\documentclass[12pt,a4paper]{article}

\usepackage[margin=2.5cm]{geometry}
\usepackage{amsmath,amssymb,amsthm,amsfonts}
\usepackage{mathtools}
\usepackage{bm}
\usepackage{xcolor}
\usepackage[colorlinks=true,linkcolor=black,citecolor=black,urlcolor=blue]{hyperref}
\usepackage{booktabs}
\usepackage{array}
\usepackage{enumitem}
\usepackage{microtype}
\usepackage{cleveref}
\usepackage{mathrsfs}
\usepackage{thmtools}
\usepackage{algorithm}
\usepackage{algpseudocode}

\theoremstyle{plain}
\newtheorem{theorem}{Theorem}[section]
\newtheorem{proposition}[theorem]{Proposition}
\newtheorem{lemma}[theorem]{Lemma}
\newtheorem{corollary}[theorem]{Corollary}

\theoremstyle{definition}
\newtheorem{definition}[theorem]{Definition}
\newtheorem{construction}[theorem]{Construction}
\newtheorem{observation}[theorem]{Observation}

\theoremstyle{remark}
\newtheorem{remark}[theorem]{Remark}

\newcommand{\GFtwo}{\mathbb{F}_2}
\newcommand{\R}{\mathbb{R}}
\newcommand{\HZ}{H_Z}
\newcommand{\HX}{H_X}
\newcommand{\rank}{\operatorname{rank}}
\newcommand{\wt}{\operatorname{wt}}
\newcommand{\dZ}{\operatorname{deg}_Z}

\title{\Large\bfseries Continuous Relaxation Decoding for CSS\\[4pt]
Quantum LDPC Codes:\\[4pt]
Energy Landscape, Basin Geometry,\\[4pt]
Gradient Dynamics, and Hardware Implementation}

\author{Yannick Schmitt \\
\texttt{ynnk.schmitt@gmail.com} \\
\vspace{4pt}
\href{https://doi.org/10.5281/zenodo.19484007}{DOI: 10.5281/zenodo.19484007}}
\date{\today}

\begin{document}
\maketitle

\begin{abstract}
We propose and rigorously analyse a continuous-relaxation decoder for CSS
quantum low-density parity-check (qLDPC) codes, with the
$[[193,25,d=4]]$ D4 causal diamond hypergraph-product code as the
primary testbed~\cite{Tillich2014,SchmittScale} and two high-distance
augmented-seed codes --- $[[112,4,(6,6)]]$ and $[[176,32,(3,6)]]$ ---
as secondary testbeds demonstrating near-term hardware relevance.
A computational characterisation of the HGP seed structure confirms
that $d=4$ is the theoretical maximum for any code built from the
causal diamond plaquette row space; a companion analysis identifies 64
weight-3 rows outside this space that raise $d_{\mathrm{cl}}$ from 4
to 6, generating the $[[112,4,(6,6)]]$ family ($kd^2/N \approx 1.29$,
the Pareto-optimal finite-block code at $N\sim100$) and the asymmetric
$[[176,32,(3,6)]]$ code (rate $0.182$, tailored to Z-biased noise).
The decoder lifts binary error variables $e_q \in \{0,1\}$ to
continuous spin variables $v_q \in \R$ and defines a squared-spring
energy $H(v)$ whose global minima are exactly the valid codewords.

Our analytical contributions are three exact theorems, each verified to
machine precision ($<10^{-10}$) across six independent CSS code families.
The critical-point theorem shows that the linear spring formulation lacks
fixed points at valid codewords, while the squared spring formulation
achieves $H(v^*)=0$ and $\nabla H(v^*)=\bm{0}$ simultaneously without
approximation.
The Hessian theorem gives the exact closed form $H_{qq}=2\lambda +
\dZ(q)/2$ at any valid codeword, implying strict positive definiteness
for all $\lambda>0$.
These results are algebraic and hold universally for any CSS code.

Our experimental contributions characterise the full geometry of the energy
landscape.
The null space of the shared-check adjacency matrix has dimension equal to
$N - \rank(\HZ)$ (109 out of 193 for the D4-HGP code), corresponding
precisely to the logical operator subspace; gradient dynamics along these
directions is governed solely by the containment term.
The direct correction path from a valid codeword to a single-error state
is provably monotone, establishing that single-error decoding failures are provably
impossible under gradient descent; empirically, this extends to $t\leq 2$ errors.
Failures for $t\geq 3$ are caused exclusively by pseudo-codeword traps
with barrier heights $\Delta E \in [1, 4]$ at $\lambda=0.5$; the actual
basin of attraction is $3$--$11\times$ wider than the Hessian radius
predicts, and both radii scale as $\lambda^{-1/2}$, making the normalised
basin shape $\lambda$-invariant.

A sweep across four optimiser variants (gradient descent, heavy-ball
momentum, Nesterov, and Langevin dynamics) reveals that heavy-ball
achieves a marginal gain at $t=5$ errors (beyond the minimum distance
$d=4$, probing optimiser behaviour under high error load) via a
path-geometry mechanism, not barrier vaulting; the Langevin decoder
requires noise temperature $T_0 \ll \Delta E_{\min} = 1.0$ for the code.
A key conceptual result is that $H(v)$ is \emph{syndrome-agnostic}: the
energy depends only on spin configuration, not on the target syndrome.
The decoder therefore acts as a \emph{codeword projector}, with its
standalone performance determined by the initialization quality rather than
by the landscape itself.

Empirically, the optimal containment coefficient follows a corrected
universal heuristic $\lambda^* \approx \tfrac{3}{16}\bar{d}$, where
$\bar{d}$ is the mean qubit degree; this refines the trace-balance
prediction $\bar{d}/4$ by a factor of $0.75$ observed consistently across
four independent code families.
Furthermore, the girth of the Tanner graph does not predict the
pseudo-codeword trap density for continuous decoding: a code with girth 8
traps at $2.5\times$ the rate of a code with girth 4, falsifying a
natural generalisation of the belief propagation analysis.

On hardware, a fully vectorised PyTorch GPU implementation achieves
$86\,535$ syndromes per second at batch size $16\,384$, a $\approx
28\,000\times$ speedup over the serial Python baseline (which operates
at $\approx 0.1\%$ theoretical FLOP efficiency due to the Python check loop).
The FLOP comparison with MWPM is regime-dependent: against the naive
$O(N^3)$ worst-case bound the continuous decoder uses $8\times$ fewer
FLOPs; against a sparse-Blossom graph model ($\approx 49\,000$ FLOPs for
810 matching-graph edges) it uses $18\times$ more.
Furthermore, standard MWPM implementations are structurally inapplicable
to the D4-HGP code: $44\%$ of its qubits have column weight $\geq 3$ and
require a hyperedge decomposition for which no simple clique approximation
is effective, yielding $\leq 31\%$ success at $t=1$ errors.
The practical step-count optimum is $K=200$ gradient steps, achieving
$96\%$ of the asymptotic success rate at twice the throughput of $K=400$.
\end{abstract}

\tableofcontents
\bigskip

\noindent\textbf{MSC 2020:} 81P70, 94B05, 90C20.\\[2pt]
\textbf{Keywords:} continuous relaxation; CSS code; qLDPC; energy landscape;
basin geometry; gradient dynamics; Hessian; pseudo-codewords;
hardware acceleration; causal diamond.

% ════════════════════════════════════════════════════════════════════════════
\section{Introduction}\label{sec:intro}
% ════════════════════════════════════════════════════════════════════════════

Decoding quantum LDPC codes is a central challenge in quantum error
correction.
The minimum-weight perfect matching (MWPM) decoder, standard for surface
codes, does not extend to the geometry of general qLDPC codes.
Belief propagation (BP) does extend, but qLDPC Tanner graphs contain
short cycles that cause BP messages to form self-reinforcing loops ---
the ``trapping set'' or ``pseudo-codeword'' phenomenon --- and BP can fail
catastrophically even at low error rates~\cite{Poulin2008,Roffe2020}.

An alternative approach, motivated by physics-based optimisation, lifts
the discrete decoding problem to a continuous energy minimisation.
Hopfield networks~\cite{Hopfield1982} encode constraint satisfaction as
energy minima; Ising machines~\cite{Inagaki2016} exploit physical dynamics
to solve combinatorial problems.
The central question for QEC is whether the continuous energy landscape can
be designed so that gradient dynamics reliably finds the correct codeword.

We apply this approach to the $[[193,25,d=4]]$ CSS qLDPC code
constructed in~\cite{SchmittScale} via the Tillich--Z\'emor hypergraph
product of the D4 causal diamond seed matrix~\cite{SchmittCSS}.
This code is a concrete, fully specified instance with known parameters.
We additionally test the decoder on two high-distance codes from the
augmented-seed family of~\cite{SchmittScale}: the $[[112,4,(6,6)]]$
code ($kd^2/N \approx 1.29$, the highest figure-of-merit code from the
causal diamond geometry at $N\sim100$) and the $[[176,32,(3,6)]]$ code
(rate $0.182$, asymmetric distances suited to Z-biased noise hardware).
For both new codes, all analytical theorems pass to machine precision,
the $\lambda^* = \tfrac{3}{16}\bar{d}$ heuristic is independently
confirmed, and the continuous decoder achieves 100\% success for
$t\leq\lfloor(d-1)/2\rfloor$ with graceful degradation beyond.
We then test the universality of all analytical results across a full
zoo of eight CSS code families including bicycle codes, repetition-code
hypergraph products, and random LDPC products.

\paragraph{Contributions.}
\begin{enumerate}[label=(\roman*),noitemsep]
\item A critical-point theorem proving that the linear spring formulation
  has no fixed points at valid codewords, while the squared spring
  formulation achieves zero energy and zero gradient simultaneously,
  without any approximation (\cref{thm:critical}).
\item An exact closed-form Hessian at any valid codeword:
  $H_{qq} = 2\lambda + \dZ(q)/2$ on the diagonal,
  implying strict local minimality for all $\lambda > 0$,
  verified to machine precision across six CSS code families
  (\cref{thm:hessian,sec:universality}).
\item An exact single-error formula showing that the gradient restoring
  force equals $\dZ(q^*)$, proven and verified to machine precision
  (\cref{thm:single_error}).
\item A three-radius characterisation of the codeword basin geometry:
  the quadratic Hessian radius $r_{\mathrm{loc}}$, the pseudo-codeword
  barrier radius $r_{\mathrm{barrier}}$, and the empirical escape radius
  $r_{\mathrm{emp}}$, with $r_{\mathrm{emp}} \gg r_{\mathrm{barrier}} \geq
  r_{\mathrm{loc}}$ and a provably monotone direct correction path for
  $t\leq 2$ (\cref{sec:basin}).
\item A comparative analysis of four optimiser variants
  (GD, heavy-ball, Nesterov, Langevin), including proof by trajectory
  analysis that the momentum mechanism is path-geometric rather than
  barrier-vaulting, and a calibration requirement for Langevin dynamics
  (\cref{sec:momentum}).
\item A conceptual identification of the decoder as a
  \emph{syndrome-agnostic codeword projector}, with implications for
  the correct measurement protocol for the decoder threshold
  (\cref{sec:threshold}).
\item A corrected universal heuristic $\lambda^* \approx \tfrac{3}{16}\bar{d}$
  for the optimal containment coefficient, and a falsification of the
  girth--trap-density hypothesis for continuous decoding (\cref{sec:universality}).
\item A hardware acceleration study demonstrating $86\,535$ syndromes/s
  on a GPU via gather/scatter vectorisation, $8\times$ fewer FLOPs than
  the naive $O(N^3)$ MWPM bound ($18\times$ \emph{more} than sparse-Blossom),
  and a $K=200$ practical optimum for step count
  (\cref{sec:hardware}).
\item Experimental characterisation of the decoder on two high-distance
  codes from the augmented-seed family~\cite{SchmittScale}: the
  $[[112,4,(6,6)]]$ and $[[176,32,(3,6)]]$ codes, both significantly
  outperforming the surface code at $N\sim100$--200.
  Decoder benchmarks independently confirm the algebraically certified
  distances and validate the $\lambda^*$ heuristic, establishing the
  continuous relaxation decoder as the natural companion algorithm for
  the Lorentzian code family (\cref{sec:universality}).
\end{enumerate}

% ════════════════════════════════════════════════════════════════════════════
\section{Setup and Notation}\label{sec:setup}
% ════════════════════════════════════════════════════════════════════════════

Let $\mathcal{C}$ be a CSS code on $N$ physical qubits with
$Z$-stabiliser matrix $\HZ \in \GFtwo^{M \times N}$ and
$X$-stabiliser matrix $\HX \in \GFtwo^{M \times N}$,
satisfying $\HZ \HX^T \equiv 0 \pmod{2}$.
We work throughout with $Z$-error decoding; the $X$-error case is symmetric.

\begin{definition}[Tanner graph and support]\label{def:tanner}
The \emph{Tanner graph} $G(\HZ)$ is bipartite with $N$ qubit nodes and
$M$ check nodes; qubit $q$ and check $j$ are adjacent iff $\HZ[j,q]=1$.
The \emph{support} of check $j$ is $\mathrm{supp}(j) = \{q : \HZ[j,q]=1\}$
with $|\mathrm{supp}(j)| = w_j$ (the check weight).
The \emph{column weight} $\dZ(q) = \sum_j \HZ[j,q]$ is the degree of
qubit $q$ in $G(\HZ)$.
The mean qubit degree is $\bar{d} = N^{-1}\sum_q \dZ(q) = M\bar{w}/N$,
where $\bar{w}$ is the mean check weight.
\end{definition}

\begin{definition}[Spin lift]\label{def:lift}
The \emph{spin representation} maps binary vectors $e \in \{0,1\}^N$ to
spin vectors $v \in \{+1,-1\}^N$ via $v_q = 1 - 2e_q$.
In this representation, the GF(2) parity constraint
$\bigoplus_{q \in \mathrm{supp}(j)} e_q = 0$
is equivalent to the multiplicative constraint
$P_j(v) := \prod_{q \in \mathrm{supp}(j)} v_q = +1$.
\end{definition}

The spin representation is exact on $\{+1,-1\}^N$; no approximation is
involved.

\begin{definition}[Shared-check adjacency matrix]\label{def:adjacency}
The \emph{shared-check adjacency matrix} $A \in \mathbb{Z}_{\geq 0}^{N\times N}$
is defined by $A_{qr} = |\{j : q\in\mathrm{supp}(j),\, r\in\mathrm{supp}(j)\}|$.
The diagonal entries satisfy $A_{qq} = \dZ(q)$.
In matrix form, $A = \HZ^T \HZ$ (as an integer matrix).
Since $A$ is a Gram matrix, $A \succeq 0$.
\end{definition}

\paragraph{The D4-HGP code.}
The primary testbed is the $[[193,25,d=4]]$ code obtained as the
hypergraph product of the $7\times 12$ independent-row basis of the D4
causal diamond $Z$-stabiliser matrix~\cite{SchmittScale}.
Parameters: $N=193$, $M=84$, $\rank(\HZ)=84$.
Column weights range from 1 to 7 with mean $\bar{d}=2.76$.
Row weights range from 5 to 11.
The null space of $A$ has dimension $\dim\ker(A) = N - \rank(\HZ) = 109$.

\begin{remark}[Corrected distance and seed characterisation]
\label{rem:d4_corrected}
An earlier version of~\cite{SchmittScale} reported $d\geq 67$ for this code.
A subsequent revision to the scaling analysis in that work established
$d=4$.
All decoder-analytic results in the present paper (Theorems~3.4, 4.1, 4.3
and all associated observations) are algebraic identities that hold for any
CSS code independent of minimum distance; they are unaffected by this
correction.
The experimental benchmarks at $t=1,2,3,4$ remain meaningful and
correctly characterise decoder performance relative to the true distance.
Experimental results at $t\geq 5$ probe the regime beyond the minimum
distance and should be interpreted as stress-tests of the optimiser, not
as below-threshold decoding.

A computational characterisation of the HGP seed structure, carried out
as part of the analysis accompanying this paper, establishes that $d=4$
is the \emph{theoretical maximum} for any hypergraph-product code built
from the causal diamond geometry.
The classical minimum distance of the primal seed
$H_Z^{\mathrm{primal}}$ (the $7\times 12$ independent-row basis of the
plaquette incidence matrix) equals $d_{\mathrm{cl}}=4$, achieved
uniquely by the three spatial-axis group vectors $G_1, G_2, G_3$; no
vector of weight $1$, $2$, or $3$ lies in $\ker(H_Z^{\mathrm{primal}})$.
Since the HGP distance satisfies
$d_Z(\mathrm{HGP}(H,H)) \leq d_{\mathrm{cl}}(H)$,
the primal seed achieves the best distance available.
The dual family of~\cite{SchmittCSS} yields only smaller classical
distances ($d_{\mathrm{cl}}=3$ for all 1248 valid Code~II X-check
matrices, verified exhaustively, and $d_{\mathrm{cl}}=2$ for the
spatial-axis seed), confirming that no orientation or variant within the
causal diamond 2-complex improves on $d=4$.
The $d=4$ result is a geometric invariant of the causal diamond, not a
defect of the Lorentzian construction.

However, a single weight-3 row \emph{outside} the plaquette row space,
whose support intersects each of $G_1, G_2, G_3$ in exactly one qubit,
kills all three minimum-weight codewords simultaneously and raises
$d_{\mathrm{cl}}$ from $4$ to $6$ (weight spectrum $\{6{:}\,12,\;
8{:}\,3\}$).
There are precisely \textbf{64} such rows of minimum weight 3.
Under the $B_4$ symmetry group (of order~96), they partition into
exactly 2 orbits: one of size~16 and one of size~48; the set is
$B_4$-closed but not $B_4$-transitive.
The hypergraph product of the resulting augmented $(8\times 12)$ seed
with itself yields a $[[208, 16]]$ code.
Because $H_{\mathrm{aug}}$ has full row rank~8, its transpose
$H_{\mathrm{aug}}^T$ has trivial kernel over $\GFtwo$, so the
Tillich--Z\'emor theorem gives $d_Z \geq \min(6,\infty)=6$ and
$d_X \geq 6$ algebraically.
A subsequent meet-in-the-middle search confirmed a weight-6
$Z$-logical operator exists, establishing $d_Z = d_X = 6$ exactly
and making this the \textbf{first code from the causal diamond geometry
exceeding the $d=4$ ceiling}: $[[208,16,6]]$ with
$k d^2/N = 16\times 36/208 \approx 2.77$.
See \cref{sec:discussion} for the full analysis and extended code families.
\end{remark}

% ════════════════════════════════════════════════════════════════════════════
\section{The Continuous Energy Function}\label{sec:energy}
% ════════════════════════════════════════════════════════════════════════════

\subsection{Two candidate energy functions}

We extend the spin representation to the full real line by introducing
continuous variables $v_q \in \R$ and defining energy functions whose
global minima are the codeword spin vectors.

\begin{definition}[Containment potential]\label{def:containment}
The \emph{containment potential} is the double-well function
\[
  V(v_q) = \frac{\lambda}{4}(v_q^2 - 1)^2, \qquad \lambda > 0,
\]
with global minima at $v_q = \pm 1$ and a local maximum at $v_q = 0$.
Its gradient is $V'(v_q) = \lambda(v_q^2-1)v_q$, which vanishes at
$v_q = \pm 1$.
\end{definition}

\begin{definition}[Linear spring energy]\label{def:linear_spring}
\[
  H_{\mathrm{lin}}(v) = \sum_{q=1}^N V(v_q) + \sum_{j=1}^M \frac{1-P_j(v)}{2}.
\]
For $v \in \{+1,-1\}^N$: each satisfied check contributes 0, each
violated check contributes 1.
\end{definition}

\begin{definition}[Squared spring energy]\label{def:sq_spring}
\[
  H(v) = \sum_{q=1}^N V(v_q) + \sum_{j=1}^M \frac{\bigl(1-P_j(v)\bigr)^2}{4}.
\]
For $v \in \{+1,-1\}^N$: each satisfied check contributes 0, each
violated check contributes 1. The discrete values of $H$ and
$H_{\mathrm{lin}}$ are identical on $\{+1,-1\}^N$.
\end{definition}

Both functions agree on the discrete constraint set.
Their critical-point structure in the continuous domain is radically different.

\subsection{Critical-point theorem}

\begin{theorem}[Critical points of the linear and squared spring energies]
\label{thm:critical}
Let $v^* \in \{+1,-1\}^N$ be a valid codeword.

\emph{(i) Linear spring:} $\nabla H_{\mathrm{lin}}(v^*) \neq \bm{0}$.
Explicitly, at $v^* = +\bm{1}$,
\[
  \frac{\partial H_{\mathrm{lin}}}{\partial v_q}\bigg|_{v=+\bm{1}}
  = -\frac{\dZ(q)}{2} \neq 0.
\]
Gradient dynamics on $H_{\mathrm{lin}}$ cannot converge to any valid codeword.

\emph{(ii) Squared spring:} $\nabla H(v^*) = \bm{0}$ and $H(v^*)=0$.
\end{theorem}

\begin{proof}
\emph{Part (i).}
The containment gradient vanishes at $v_q=\pm 1$: $V'(\pm 1) = 0$.
The linear spring gradient at qubit $q$ evaluated at $v=+\bm{1}$ is
\[
  -\frac{1}{2} \sum_{j: q\in\mathrm{supp}(j)} \prod_{k \in \mathrm{supp}(j),\, k\neq q} v_k
  = -\frac{\dZ(q)}{2},
\]
since every $v_k = +1$.

\emph{Part (ii).}
The squared spring gradient at qubit $q$ is
\[
  \frac{\partial H}{\partial v_q} = V'(v_q) -
    \frac{1}{2}\sum_{j: q\in\mathrm{supp}(j)}
    \bigl(1 - P_j\bigr)\cdot \frac{P_j}{v_q}.
\]
At a valid codeword $P_j = +1$ for all $j$, so $(1-P_j) = 0$, and
the spring gradient vanishes.
$V'(v_q^*) = 0$ since $v_q^* = \pm 1$.
Thus $\nabla H(v^*) = \bm{0}$.
$H(v^*)=0$ follows from $V(\pm 1)=0$ and $(1-P_j)^2/4 = 0$.
\end{proof}

\begin{remark}[Exact reformulation, not approximation]
The squared spring is not an approximation of the linear spring.
Both equal the check-violation count on $\{+1,-1\}^N$.
The term $(1-P_j)^2/4$ equals $(1-\cos\theta_j)/2$ where
$\theta_j \in \{0,\pi\}$ is the phase of $P_j$ on $\{+1,-1\}^N$,
the XZ-rotation energy familiar from quantum Ising models.
\end{remark}

% ════════════════════════════════════════════════════════════════════════════
\section{Analytical Results}\label{sec:analysis}
% ════════════════════════════════════════════════════════════════════════════

\subsection{Single-error energy and gradient}

\begin{theorem}[Single-error formulae]\label{thm:single_error}
Let $v^{(q^*)}$ denote the spin vector with $v_{q^*} = -1$ and
$v_q = +1$ for $q \neq q^*$. Then:
\begin{align}
  H(v^{(q^*)}) &= \dZ(q^*), \label{eq:single_energy} \\
  \frac{\partial H}{\partial v_q}\bigg|_{v = v^{(q^*)}} &=
    \begin{cases}
      -\dZ(q^*) & q = q^*, \\
      +A_{q,q^*} & q \neq q^*, \\
    \end{cases} \label{eq:single_grad}
\end{align}
\end{theorem}

\begin{proof}
\emph{Energy.}
Each check $j$ containing $q^*$ has $P_j = -1$,
contributing $(1-(-1))^2/4 = 1$.
Checks not containing $q^*$ have $P_j = +1$, contributing 0.
The containment contributes $V(-1) = 0$.
Thus $H(v^{(q^*)}) = \dZ(q^*)$.

\emph{Gradient at $q = q^*$.}
For each $j \ni q^*$: $P_j = -1$,
$\frac{\partial}{\partial v_{q^*}} \frac{(1-P_j)^2}{4}
  = -\frac{1}{2}(1-P_j)\frac{P_j}{v_{q^*}} = -1$.
Summing over $\dZ(q^*)$ checks gives $-\dZ(q^*)$.

\emph{Gradient at $q \neq q^*$.}
For $j \ni q, q^*$: $P_j = -1$ (since $v_{q^*}=-1$ and all others $+1$)
and $\frac{P_j}{v_q} = -1/1 = -1$,
contributing $-\frac{1}{2}(1-(-1))(-1) = -\frac{1}{2}(2)(-1) = +1$.
For $j \ni q, j \not\ni q^*$: $P_j = +1$, contributing 0.
Summing over the $A_{q,q^*} = |\{j : q,q^* \in \mathrm{supp}(j)\}|$ shared
checks gives $+A_{q,q^*}$.
\end{proof}

\begin{corollary}[Restoring force magnitude]
The restoring force on the erroneous qubit is $+\dZ(q^*) > 0$,
directed toward $v_{q^*} = +1$.
Its magnitude equals the qubit's Tanner graph degree.
The gradient at a neighbouring qubit $q \neq q^*$ is $+A_{q,q^*} \geq 0$:
the spring term \emph{pushes} neighbours away from $+1$ (toward the error),
consistent with the squared-spring trying to restore violated checks by
flipping all participants.
\end{corollary}

\subsection{Hessian at valid codewords}

\begin{theorem}[Exact Hessian at any valid codeword]\label{thm:hessian}
At any valid codeword $v^*$, the Hessian $\mathbf{H} = \nabla^2 H(v^*)$
has the exact structure:
\begin{align}
  \mathbf{H}_{qq} &= 2\lambda + \frac{\dZ(q)}{2}, \label{eq:hess_diag} \\
  \mathbf{H}_{qr} &= \frac{A_{qr}}{2}, \quad q \neq r,
    \label{eq:hess_offdiag}
\end{align}
where $A = \HZ^T\HZ$ is the shared-check adjacency matrix.
In matrix form: $\mathbf{H} = 2\lambda I + \tfrac{1}{2}A$.
\end{theorem}

\begin{proof}
\emph{Diagonal.}
The containment contributes $V''(\pm 1) = \lambda(3\cdot 1-1) = 2\lambda$.
For the spring term, at $P_j = +1$:
$\frac{\partial^2}{\partial v_q^2}\frac{(1-P_j)^2}{4} = \frac{1}{2}\left(\frac{P_j}{v_q}\right)^2 = \frac{1}{2}$.
Summing over $\dZ(q)$ checks gives $\dZ(q)/2$.

\emph{Off-diagonal.}
Only checks containing both $q$ and $r$ contribute.
At $P_j = +1$:
$\frac{\partial^2}{\partial v_q \partial v_r}\frac{(1-P_j)^2}{4}
  = \frac{1}{2}\cdot \frac{P_j^2}{v_q v_r} = \frac{1}{2}$,
since $v_q^*,v_r^* = \pm 1$ and $P_j^2 = 1$.
Summing over $A_{qr}$ shared checks gives $A_{qr}/2$.
\end{proof}

\begin{corollary}[Strict local minimality for all $\lambda > 0$]\label{cor:local_min}
$\mathbf{H} = 2\lambda I + \tfrac{1}{2}A \succ 0$ for all $\lambda > 0$,
since $A = \HZ^T\HZ \succeq 0$.
Every valid codeword is a strict local minimum of $H(v)$.
\end{corollary}

\subsection{Hessian spectral structure}

The eigenvalue formula $\mu_k(\lambda) = 2\lambda + \sigma_k(A/2)$ follows
directly from \cref{thm:hessian}, where $\sigma_k(A/2)$ are the eigenvalues
of $A/2$.
These are code-intrinsic ($\lambda$-independent) and can be computed once
from $\HZ$.

\begin{proposition}[Null-space structure]\label{prop:null}
The null space of $A = \HZ^T\HZ$ over $\R$ has dimension
$\dim\ker(A) = N - \rank_{\R}(\HZ)$,
and corresponds precisely to vectors that share no check pair.
For the D4-HGP code: $\dim\ker(A) = 193 - 84 = 109$.
In the null directions, $\mu_{\mathrm{null}} = 2\lambda$ exactly, so
basin curvature is governed solely by the containment term.
These null directions span a subspace that contains the logical operator
representatives.
\end{proposition}

\begin{proof}
$\ker(A) = \ker(\HZ)$ over $\R$ since $A = \HZ^T\HZ$ and
$Ax=0 \Leftrightarrow x^T A x = \|\HZ x\|^2 = 0 \Leftrightarrow \HZ x = 0$.
The dimension of $\ker(\HZ)$ over $\R$ is $N - \rank_{\R}(\HZ)$.
Note that for $\{0,1\}$-matrices arising from CSS codes, $\rank_{\R}(\HZ)
\geq \rank_{\GFtwo}(\HZ)$; the two ranks coincide whenever the real row
space contains no non-trivial integer linear dependencies beyond those
detectable over $\GFtwo$.
For all codes in this paper (verified numerically), $\rank_{\R} = \rank_{\GFtwo}$.
The containment contribution to eigenvalue in these directions is $2\lambda$;
the spring contribution is $x^T (A/2) x = 0$.
\end{proof}

\begin{remark}[Code structure as geometric advantage]
For the D4-HGP code, $109/193 \approx 56\%$ of Hessian eigendirections
lie in the null space of $A$.
In these directions the codeword basin extends to the containment clip
boundary, measured empirically as $r_{\mathrm{emp}} \approx \rho/\sqrt{\lambda}$
where $\rho = 1.5$ is the clip parameter.
This is the code's topological structure acting as a structural advantage
for the continuous decoder.
\end{remark}

\subsection{Local minima and pseudo-codewords}

\begin{proposition}[Local minima are pseudo-codewords]\label{prop:pseudo}
Let $v^{\mathrm{loc}}$ be a local minimum of $H$ and let
$\hat{v} = \mathrm{sign}(v^{\mathrm{loc}})$.
If $\hat{v}$ is not a valid codeword, then the binary vector
$\hat{e} = (1-\hat{v})/2 \bmod 2$ is a pseudo-codeword:
$\HZ\hat{e} \not\equiv \bm{0}$ (mod 2).
\end{proposition}

This is the continuous analogue of BP trapping sets.
Unlike BP fixed points, these traps are geometric objects (local basins)
amenable to physical escape strategies.

% ════════════════════════════════════════════════════════════════════════════
\section{The Gradient Descent Decoder}\label{sec:decoder}
% ════════════════════════════════════════════════════════════════════════════

\begin{construction}[Continuous Relaxation Decoder]\label{con:decoder}
Given syndrome $s \in \{0,1\}^M$:
\begin{enumerate}[noitemsep]
\item \emph{Initialise:} $v^{(0)} = +\bm{1} - 2\hat{e}^{(0)}$, where $\hat{e}^{(0)}$
  is a syndrome-consistent error estimate (e.g.\ greedy peeling, see \cref{sec:threshold}).
\item \emph{Iterate:}
  $v^{(k+1)} = \mathrm{clip}\bigl(v^{(k)} - \eta \nabla H(v^{(k)}),\;
  [-\rho, \rho]\bigr)$,
  with $\eta=0.05$, $\rho=1.5$, for $K$ steps.
\item \emph{Snap:} $\hat{v} = \mathrm{sign}(v^{(K)})$;
  $\hat{e} = (1-\hat{v})/2 \bmod 2$.
\item Apply $\hat{e}$ as the correction.
\end{enumerate}
\end{construction}

\paragraph{Gradient computation.}
The gradient $\nabla H$ is computed analytically from the squared-spring
formula.
The partial derivative at qubit $q$ is
\begin{equation}
  \frac{\partial H}{\partial v_q}
  = \lambda(v_q^2-1)v_q - \frac{1}{2}
    \sum_{j\,:\,q\in\mathrm{supp}(j)} (1-P_j)\cdot\frac{P_j}{v_q}.
  \label{eq:grad_formula}
\end{equation}
No message-passing or graph traversal is needed.

\paragraph{Complexity.}
Each gradient step requires $3N$ FLOPs for the containment term and
$3\cdot\mathrm{nnz}(\HZ)$ FLOPs for the spring term, giving
$3(N + \mathrm{nnz}(\HZ))$ FLOPs per step.
For the D4-HGP code ($N=193$, $\mathrm{nnz}=532$):
$3\times 725 = 2{,}175$ FLOPs per step.
At $K=400$ steps: $870{,}000$ FLOPs per syndrome,
versus $O(N^3) \approx 7.2\times10^6$ for MWPM.

% ════════════════════════════════════════════════════════════════════════════
\section{Basin Geometry and Convergence Theory}\label{sec:basin}
% ════════════════════════════════════════════════════════════════════════════

This section characterises the geometry of the codeword basin of attraction
through three complementary radii and proves that decoding failures are
impossible for $t \leq 2$ errors under gradient descent.

\subsection{Three nested radii}

\begin{definition}[Three convergence radii]\label{def:radii}
Fix a valid codeword $v^*$ and $\lambda > 0$. Define:
\begin{itemize}[noitemsep]
\item $r_{\mathrm{loc}}(\lambda) = \sqrt{2/\mu_{\min}(\lambda)}$, the
  \emph{Hessian radius}, where $\mu_{\min} = 2\lambda + \sigma_{\min}(A/2)$
  is the smallest Hessian eigenvalue.
\item $r_{\mathrm{barrier}}(\lambda) = \sqrt{2\Delta E / \mu_{\min}(\lambda)}$,
  the \emph{barrier radius} corresponding to saddle height $\Delta E$ to
  the nearest pseudo-codeword basin.
\item $r_{\mathrm{emp}}(\lambda, \hat{u})$, the \emph{empirical escape radius}
  along direction $\hat{u}$: the smallest $r$ such that
  $\mathrm{GD}(v^* + r\hat{u})$ fails to return to any valid codeword.
\end{itemize}
\end{definition}

\begin{observation}[Nested basin geometry]\label{obs:nested}
For the D4-HGP code at $\lambda = 0.5$:
\[
  r_{\mathrm{loc}} \approx 1.41
  \;\leq\;
  r_{\mathrm{barrier}} \in [1.41,\, 2.83]
  \;\ll\;
  r_{\mathrm{emp}} \in [7,\, 15] \times r_{\mathrm{loc}},
\]
across $\hat{u}$ ranging from the stiffest to the softest Hessian
eigendirection.
The Hessian radius is a conservative lower bound; the actual basin is
$3$--$11\times$ wider than $r_{\mathrm{loc}}$.
\end{observation}

The numerical values are reported for 12 sampled pseudo-codewords at
$\lambda=0.5$ in \cref{tab:barriers}.

\subsection{Monotone direct correction path}

\begin{proposition}[No barrier on the direct correction path]\label{prop:monotone}
Consider the one-dimensional restriction of $H$ to the path
$v^*(t) = v^* + t(v^{(q^*)} - v^*)$ for $t\in[0,1]$,
connecting a valid codeword $v^*$ to a single-error state $v^{(q^*)}$.
For any physical parameter values satisfying $\lambda < 2\dZ(q^*)$,
the gradient $dH/dt$ is everywhere negative on $(0,1)$: the path is
\emph{strictly monotone decreasing} in energy from $v^{(q^*)}$ to $v^*$.
In particular, no saddle point separates $v^{(q^*)}$ from $v^*$ along
this path.
\end{proposition}

\begin{proof}
Restricting to the 1D path by holding all other qubits at $+1$, the
energy along the path is
\[
  E(s) = \frac{\lambda}{4}(s^2-1)^2 + \frac{\dZ(q^*)}{4}(1-s)^2,
  \quad s \in [-1, +1],
\]
where $s = v_{q^*}(t) = -1 + 2t$ parametrises the erroneous qubit.
A stationary point other than $s=+1$ (the codeword) satisfies the
cubic equation $(s-1)[\lambda s(s+1) + \dZ(q^*)/2] = 0$.
The bracket has discriminant $\Delta = \lambda^2 - 2\lambda\cdot\dZ(q^*)$.
For $\lambda < 2\dZ(q^*)$, which holds for all physical $\lambda$ since
$\lambda \ll \dZ(q^*)$ in the interesting regime, we have $\Delta < 0$:
the bracket has no real roots other than the trivial factored root.
Therefore $s=+1$ is the unique stationary point on $[-1,+1]$.
\end{proof}

\begin{corollary}[Provably perfect decoding for $t=1$; empirical for $t=2$]\label{cor:perfect}
For a code with minimum qubit degree $d_{\min} \geq 1$, gradient descent
with any $\lambda < 2d_{\min}$ corrects every single-error pattern at
100\% rate.
Empirically, perfect decoding extends to $t=2$ errors in all experiments;
however, the monotone-path argument does not constitute a formal proof
for $t=2$, since correcting two simultaneous errors requires showing that
the joint two-error initial point lies in the codeword basin, not merely
that each individual sub-problem is monotone.
\end{corollary}

This explains the experimental observation of 100\% success for $t\leq 2$
errors over all 30 random trials at each weight.

\subsection{Pseudo-codeword barriers}

Failures for $t\geq 3$ arise not from the direct correction path but from
the trajectory landing in the basin of a pseudo-codeword $v_{\mathrm{pc}}$
whose snapped configuration violates at least one check.

\begin{observation}[Barrier height distribution]\label{obs:barriers}
For 12 sampled pseudo-codewords collected from failed $t=5$ trials at
$\lambda=0.5$, the straight-line energy maximum $\Delta E = \max_{t\in[0,1]}
H(v^* + t(v_{\mathrm{pc}} - v^*))$ satisfies:
$\Delta E \in [1.0, 4.0]$, with mean $1.67$.
The tightest trap has $\Delta E = 1.0$ and involves a 2-qubit pattern
sharing a single Tanner graph cycle of length 4 (girth cycle).
The corresponding barrier radius $r_{\mathrm{barrier}} = r_{\mathrm{loc}} = 1.41$,
confirming that the Hessian radius exactly predicts the binding constraint
for the critical trap class.
\end{observation}

\begin{remark}[Barrier measurement note]
The reported $\Delta E$ values are upper bounds on the true saddle height,
since the straight-line path is not guaranteed to be the minimum energy
path between $v^*$ and $v_{\mathrm{pc}}$.
The true (minimum-energy-path) saddle may be lower.
An exact saddle search would require solving for the index-1 critical
point of $H$ between the two basins.
\end{remark}

\subsection{Scaling with $\lambda$}

\begin{observation}[$\lambda$-invariance of normalised basin shape]
\label{obs:scaling}
For both $r_{\mathrm{loc}}$ and $r_{\mathrm{barrier}}$, the scaling law
$r \propto \lambda^{-1/2}$ holds across $\lambda \in [0.1, 2.0]$.
The ratio $r_{\mathrm{barrier}} / r_{\mathrm{loc}}$ is nearly constant
over this range ($\lesssim 10\%$ variation).
Consequently, the basin geometry---when expressed in units of
$r_{\mathrm{loc}}$---is $\lambda$-invariant.
This implies that optimising $\lambda$ does not fundamentally change the
barrier-to-basin-size ratio; it only scales all length scales together.
\end{observation}

\begin{table}[ht]
\centering
\caption{Basin geometry parameters at $\lambda = 0.5$ for the D4-HGP code.
  Barrier heights measured along the straight-line path $v^* \to \mathrm{snap}(v_{\mathrm{pc}})$.
  The Hessian radius $r_{\mathrm{loc}} = \sqrt{2/\mu_{\min}} = 1.414$ at $\lambda=0.5$.}
\label{tab:barriers}
\small
\begin{tabular}{cccccc}
\toprule
PC type & $n_{\mathrm{flip}}$ & $\Delta E$ & $r_{\mathrm{barrier}}$ & $r_{\mathrm{barrier}}/r_{\mathrm{loc}}$ \\
\midrule
2-flip (most common) & 2 & 1.00 & 1.41 & 1.00 \\
4-flip, interior saddle & 4 & $2.00$--$2.27$ & $2.00$--$2.13$ & $1.41$--$1.51$ \\
6-flip, interior saddle & 6 & 2.51 & 2.24 & 1.58 \\
4-flip, max syndrome wt & 4 & 4.00 & 2.83 & 2.00 \\
\midrule
\multicolumn{2}{l}{Summary} & min 1.00, max 4.00 & min 1.41, max 2.83 & 1.00--2.00 \\
\bottomrule
\end{tabular}
\end{table}

\begin{table}[ht]
\centering
\caption{Hessian spectrum and quadratic convergence radius at selected $\lambda$.
  $\sigma_{\min}(A/2) \approx 0$ (null space dim = 109);
  $r_{\mathrm{loc}} \approx 1/\sqrt{\lambda}$ for all physical $\lambda$.}
\label{tab:spectrum}
\small
\begin{tabular}{ccccc}
\toprule
$\lambda$ & $\mu_{\min}$ & $\mu_{\max}$ & $\kappa$ & $r_{\mathrm{loc}}$ \\
\midrule
0.05 & 0.100 & 14.99 & 149.9 & 4.472 \\
0.10 & 0.200 & 15.09 & 75.5 & 3.162 \\
0.50 & 1.000 & 15.89 & 15.89 & 1.414 \\
1.00 & 2.000 & 16.89 & 8.44 & 1.000 \\
2.00 & 4.000 & 18.89 & 4.72 & 0.707 \\
\bottomrule
\end{tabular}
\end{table}

% ════════════════════════════════════════════════════════════════════════════
\section{Experiments on the $[[193,25]]$ Code}\label{sec:experiments}
% ════════════════════════════════════════════════════════════════════════════

\subsection{Setup}

All experiments use $\eta = 0.05$, $K = 400$ steps, $\rho = 1.5$, $\lambda = 0.5$
unless otherwise stated.
Error patterns are drawn uniformly at random.
Trial counts are specified per experiment.

\subsection{Analytical predictions verified}

\cref{tab:analytical} reports the numerical verification of all analytical
results from \cref{sec:analysis}.
All results match to machine precision.

\begin{table}[ht]
\centering
\caption{Numerical verification of analytical formulae on the $[[193,25]]$ code.
  All discrepancies are below machine epsilon ($<10^{-10}$).}
\label{tab:analytical}
\small
\begin{tabular}{lcc}
\toprule
Claim & Predicted & Numerical \\
\midrule
Ground state energy $H(+\bm{1})$ & $0$ & $0$ (exact) \\
Ground state gradient $\|\nabla H(+\bm{1})\|$ (squared spring) & $0$ & $1.76\times10^{-13}$ \\
Ground state gradient $\|\nabla H_{\mathrm{lin}}(+\bm{1})\|$ & $\sqrt{\sum_q d_q^2}/2$ & $22.54$ (matches) \\
Single-error energy, $d_q = 7$ & $7$ & $7.000$ (exact) \\
Single-error energy, $d_q = 1$ & $1$ & $1.000$ (exact) \\
Single-error gradient at $q^*$ ($d_q=7$) & $-7$ & $-7.000$ (exact) \\
Hessian diagonal ($\lambda=1$, $d_q=7$) & $5.5$ & $5.500$ (exact) \\
Hessian diagonal ($\lambda=1$, $d_q=1$) & $2.5$ & $2.500$ (exact) \\
\bottomrule
\end{tabular}
\end{table}

\subsection{Gradient descent convergence vs.\ error weight}

\begin{table}[ht]
\centering
\caption{GD decoder performance vs.\ error weight $t$.
  $\lambda=0.5$, $K=400$, 100 trials each.}
\label{tab:convergence}
\small
\begin{tabular}{cccc}
\toprule
$t$ & Success rate & Mean residual syndrome (failures) & Initial energy \\
\midrule
1 & 100/100 (1.00) & -- & $\approx 2.64$ \\
2 & 100/100 (1.00) & -- & $\approx 4.68$ \\
3 & 99/100 (0.99) & 0.14 & $\approx 7.40$ \\
4 & 91/100 (0.91) & 0.18 & $\approx 10.1$ \\
5 & 82/100 (0.82) & 0.08 & $\approx 12.84$ \\
\bottomrule
\end{tabular}
\end{table}

For $t\leq 2$ the decoder is exact over all 100 trials, consistent with
\cref{cor:perfect}.
Failures at $t\geq 3$ correspond exclusively to pseudo-codeword trapping
(nonzero residual syndrome).

\subsection{Lambda sensitivity}\label{sec:lambda}

\cref{tab:lambda} reports success rate vs.\ $\lambda$ for $t=3$ errors.

\begin{table}[ht]
\centering
\caption{Decoder success rate vs.\ $\lambda$ for $t=3$ errors (100 trials).
  The heuristic optimal $\lambda^* \approx \tfrac{3}{16}\bar{d} = 0.517$.}
\label{tab:lambda}
\small
\begin{tabular}{ccc}
\toprule
$\lambda$ & Success rate & Mean final energy \\
\midrule
0.10 & 97/100 (0.97) & 0.017 \\
0.50 & 99/100 (0.99) & 0.000 \\
1.00 & 78/100 (0.78) & 0.311 \\
2.00 & 32/100 (0.32) & 0.775 \\
5.00 & 8/100 (0.08) & 2.607 \\
\bottomrule
\end{tabular}
\end{table}

The optimal range is $\lambda \in [0.1, 0.5]$.
The mechanistic interpretation follows from \cref{thm:hessian}: at large
$\lambda$ the deep basin curvature $2\lambda + \dZ(q)/2$ immediately
suppresses spring-driven corrections.
A heuristic balance condition $2\lambda = \bar{d}/2$ yields
$\lambda_{\mathrm{pred}} = \bar{d}/4 = 0.69$ for the D4-HGP code.
Empirically the optimum sits at $\lambda^* \approx 0.75 \times 0.69 = 0.52$
(see \cref{sec:universality} for the full derivation of this factor).

% ════════════════════════════════════════════════════════════════════════════
\section{Momentum-Based Decoding}\label{sec:momentum}
% ════════════════════════════════════════════════════════════════════════════

\subsection{Optimiser variants}

Three momentum-based optimisers are compared against vanilla GD.

\paragraph{Heavy-ball (HB; Polyak 1964).}
$v^{(k+1)} = v^{(k)} - \eta\nabla H(v^{(k)}) + \gamma(v^{(k)} - v^{(k-1)})$.
The momentum term $\gamma(v^{(k)}-v^{(k-1)})$ injects velocity proportional
to the current velocity.
For a quadratic landscape with condition number $\kappa$, the optimal
parameter is $\gamma_{\mathrm{opt}} = \left(\frac{\sqrt{\kappa}-1}{\sqrt{\kappa}+1}\right)^2$.
For the D4-HGP code at $\lambda=0.5$, $\kappa = 15.89$, giving
$\gamma_{\mathrm{opt}} \approx 0.35$.
Periodic momentum restarts (every $r$ steps) prevent orbit trapping.

\paragraph{Nesterov accelerated gradient (NAG).}
$y^{(k)} = v^{(k)} + \frac{k-1}{k+2}(v^{(k)}-v^{(k-1)})$, then
$v^{(k+1)} = y^{(k)} - \eta\nabla H(y^{(k)})$,
with adaptive restart on energy increase.
Achieves $O(1/k^2)$ convergence in quadratic landscapes.

\paragraph{Annealed Langevin dynamics.}
$v^{(k+1)} = v^{(k)} - \eta\nabla H(v^{(k)}) + \sqrt{2\eta T_k}\,\xi^{(k)}$,
$\xi^{(k)}\sim\mathcal{N}(0,I)$,
with exponential annealing $T_k = T_0(T_{\infty}/T_0)^{k/K}$.
The Kramers rate for barrier crossing is $\propto \exp(-\Delta E/T)$.

\subsection{Main comparison}

\begin{table}[ht]
\centering
\caption{Decoder success rates (100 trials each, $\lambda=0.5$).
  HB parameters: $\gamma=0.85$, restart every 150 steps.
  Langevin parameters: $T_0=0.5$, $T_\infty=0.01$, 800 steps.}
\label{tab:momentum}
\small
\begin{tabular}{lcccc}
\toprule
Method & $t=3$ & $t=4$ & $t=5$ & Mean (3--5) \\
\midrule
GD (baseline) & 0.990 & 0.910 & 0.820 & 0.907 \\
Heavy-ball ($\gamma=0.85$) & 0.980 & 0.910 & 0.850 & 0.913 \\
Nesterov (adaptive) & 0.980 & 0.900 & 0.830 & 0.903 \\
Langevin ($T_0=0.5$) & 0.100 & 0.090 & 0.140 & 0.110 \\
\midrule
\multicolumn{4}{l}{All decoders (except Langevin) achieve 100/100 at $t=1,2$} \\
Langevin correctness check & 9/30 & 2/30 & -- & -- \\
\bottomrule
\end{tabular}
\end{table}

Heavy-ball achieves a $+3\%$ improvement over GD at $t=5$ (0.850 vs.\
0.820).
The magnitude corresponds to approximately $0.6–0.8\sigma$ at 100 trials
($\sigma \approx \sqrt{0.85\times 0.15/100} \approx 3.6\%$), so this
result requires confirmation with $N\geq 500$ trials.

\textbf{Langevin failure.}
At $T_0 = 0.5$, the Langevin decoder fails the $t=1$ correctness check
catastrophically (9/30), demonstrating that noise above
$\Delta E_{\min} = 1.0$ destroys easily correctable patterns.
The required calibration condition is $T_0 \ll \Delta E_{\min}$, derived
analytically from \cref{tab:barriers}: $T_0 \ll 1.0$, so $T_0 \leq 0.1$.
At $T_0 = 0.05$ (near-zero noise, approaching GD), Langevin achieves
$88\%$ at $t=4$ --- consistent with the GD baseline.

\subsection{The path-geometry mechanism}

\begin{observation}[No barrier vaulting]\label{obs:no_vault}
For a representative $t=5$ trial where GD fails but heavy-ball succeeds,
the energy trajectory satisfies:
\begin{align*}
  E_{\mathrm{start}} &= 9.000, \\
  E_{\mathrm{peak,\,HB}} &= 9.000 \quad \text{(overshoot } = 0.000\text{)}, \\
  E_{\mathrm{final,\,HB}} &= 0.000 \quad \text{(convergence by step 100)}.
\end{align*}
The peak energy never exceeds the starting energy.
Heavy-ball does \emph{not} escape pseudo-codeword traps by kinetic
overshoot over the energy barrier; instead, it follows a qualitatively
different gradient trajectory that avoids entering the trap basin.
\end{observation}

This falsifies the intuitive ``kinetic barrier crossing'' hypothesis for
the heavy-ball mechanism in this landscape.
The physical mechanism---whether momentum prevents the trajectory from
making sharp high-curvature turns into narrow pseudo-codeword basins, or
whether it takes a qualitatively different path that bypasses the basin
entirely---remains an open problem.

\subsection{Parameter sensitivity}

\begin{observation}[Flat $\gamma$ response]\label{obs:gamma_flat}
The heavy-ball success rate at $t=4$ is constant at $\approx 0.933$ for
$\gamma \in [0.50, 0.85]$, then declines sharply above $\gamma = 0.90$.
The theoretical optimum $\gamma_{\mathrm{opt}} \approx 0.35$
(from $\kappa = 15.89$) lies far outside the effective regime.
The quadratic $\kappa$-based formula is inapplicable to the nonlinear
spring-product landscape.
\end{observation}

\begin{observation}[Langevin noise monotonicity]\label{obs:langevin_mono}
The Langevin success rate at $t=4$ decreases monotonically with $T_0$
across $T_0 \in [0.05, 1.0]$ (from 0.883 to 0.133).
More noise is always worse.
No Kramers-style barrier crossing was observed: the stochastic escape rate
did not produce a non-monotone benefit at intermediate temperatures.
\end{observation}

\begin{observation}[Restart insensitivity]\label{obs:restart}
The heavy-ball success rate at $t=5$ is identically $0.800 \pm 0.004$
for all restart intervals tested (none, 50, 100, 150, 200, 300 steps).
Periodic restarts have no measurable effect on this landscape.
\end{observation}

% ════════════════════════════════════════════════════════════════════════════
\section{Initialization, Threshold, and Failure Modes}\label{sec:threshold}
% ════════════════════════════════════════════════════════════════════════════

\subsection{The syndrome-agnostic landscape and the projector paradigm}

\begin{observation}[Syndrome-agnostic energy]\label{obs:agnostic}
The energy function $H(v)$ depends only on the spin configuration $v$
and not on the target syndrome $s$:
\[
  H(v) = \sum_q \frac{\lambda}{4}(v_q^2-1)^2 + \sum_j \frac{(1-P_j)^2}{4}.
\]
The syndrome $s$ enters only in the initialisation $v^{(0)}$, not in the
gradient dynamics.
If the initialisation is at a global minimum ($v^{(0)} = \pm\bm{1}^N$,
a valid codeword), the gradient is exactly zero and the decoder does not
move.
\end{observation}

This has a fundamental implication: the decoder does not ``search'' for
a syndrome-consistent solution by gradient dynamics.
It \emph{projects} the initial configuration into the nearest codeword
basin.
We call this the \emph{projector paradigm}.

\begin{definition}[Decoder Saturation Point (DSP)]
The \emph{decoder saturation point} $p_{\mathrm{DSP}}$ is the physical
error rate $p$ at which the total logical failure rate $P_L(p) = 0.50$
for a fixed code.

\emph{Note on terminology.}
The term ``pseudo-threshold'' is reserved in the QEC literature for the
crossing $P_L(p_{\mathrm{pseudo}}) = p$ between the logical and physical
error rates.
This is distinct from $p_{\mathrm{DSP}}$ and requires a code family
(varying $N$), not a single code.
$p_{\mathrm{DSP}}$ is a single-code engineering figure of merit.
\end{definition}

\begin{observation}[DSP measurement]
For the D4-HGP code under independent $Z$-noise with greedy peeling
initialisation and vanilla GD decoding (200 trials per $p$), the failure
rate $P_L(p)$ is monotone increasing.
The DSP satisfies $p_{\mathrm{DSP}} \lesssim 0.05$, providing an upper
bound on the true threshold $p_c$, which requires a code family with
growing $N$.
\end{observation}

\subsection{Greedy peeling initialisation}

\begin{construction}[Greedy peeling]\label{con:peeling}
Given syndrome $s\in\{0,1\}^M$:
\begin{enumerate}[noitemsep]
\item Initialise $\hat{e} = \bm{0}$.
\item Repeat over unsatisfied checks $j$ where $(\HZ\hat{e})_j \neq s_j$:
  select the lowest-degree qubit $q$ in $\mathrm{supp}(j)$, flip $\hat{e}_q$.
  Update residual $r = (s \oplus \HZ\hat{e}) \bmod 2$.
\item If any checks remain unsatisfied after two passes, fall back to
  $\hat{e} = \bm{0}$ (valid-codeword fallback).
\end{enumerate}
\end{construction}

The low-degree preference minimises collateral check violations.
Validated to produce a syndrome-consistent guess in $\geq 90\%$ of
random trials at $p \leq 0.02$.
When peeling fails and the fallback $\hat{e}=\bm{0}$ is used, the
decoder starts at a global minimum and cannot correct the error.

\subsection{Failure mode census}

\begin{observation}[Dominant failure modes by error rate]\label{obs:failures}
At 300 trials per $p$:
\begin{itemize}[noitemsep]
\item At $p=0.03$: syndrome failures (GD trapped at pseudo-codeword) dominate.
  Logical-class failures are rare.
\item At $p=0.07$: logical failures (wrong logical class despite satisfied
  syndrome) become significant, consistent with the error weight exceeding
  the effective correction radius.
\end{itemize}
The crossover from ``syndrome-failure-dominated'' to
``logical-failure-dominated'' occurs near $p\approx 0.05$.
\end{observation}

\subsection{Algorithmic vs.\ landscape failures}

\begin{observation}[n-steps classification]\label{obs:nsteps}
Sweeping step count $K \in \{100, 200, 400, 800\}$ at fixed $p$:
\begin{itemize}[noitemsep]
\item At $p \leq 0.03$: $P_L$ decreases substantially with $K$ (from
  $\sim 0.40$ at $K=100$ to $\sim 0.15$ at $K=800$).
  Failures are \emph{algorithmic}: the landscape has the correct minimum,
  but the optimiser exits early.
\item At $p \geq 0.07$: $P_L$ is nearly flat across $K$.
  Failures are \emph{landscape-determined}: logical errors are the
  global minima and no amount of additional optimisation helps.
\end{itemize}
Crossings of $P_L(p)$ curves for different $K$ values reflect algorithmic
capacity limits, not finite-size-scaling threshold physics.
The true threshold $p_c$ requires a code family with growing $N$.
\end{observation}

\subsection{Path forward: BP-initialised continuous decoder}

Because $H(v)$ is a projector, the appropriate use of the continuous
decoder is as a post-processor for belief propagation, analogous to
BP-OSD~\cite{Roffe2020}.
BP provides a fractional confidence vector $v^{(0)} \in (-1,+1)^N$ already
biased toward the correct basin; the continuous decoder then applies a
single smooth gradient flow to the nearest codeword.
Measuring $P_L(p)$ for this combined decoder (and locating the crossing
point across a code family) is the correct protocol for determining $p_c$.

% ════════════════════════════════════════════════════════════════════════════
\section{CSS Code Universality}\label{sec:universality}
% ════════════════════════════════════════════════════════════════════════════

\subsection{Code zoo}

\cref{tab:zoo} describes the eight code families used for universality testing,
including the two new high-distance augmented-seed codes.

\begin{table}[ht]
\centering
\caption{CSS code zoo used for universality verification.
  $\bar{d}$: mean qubit degree in $\HZ$.
  ``Bicycle-$n$ ($k=0$)'' indicates rank-deficient codes included as
  algebraic testbeds; their decoding success rates are not physically
  meaningful.
  \textbf{Bold rows}: new high-distance augmented-seed codes from~\cite{SchmittScale}.}
\label{tab:zoo}
\small
\begin{tabular}{lccccccc}
\toprule
Code & $N$ & $M$ & $k$ & $\bar{d}$ & $\lambda^*_{\mathrm{pred}}$ & $\dim\ker(A)$ & $\kappa(\lambda^*)$ \\
\midrule
D4-HGP $[[193,25]]$ & 193 & 84 & 25 & 2.76 & 0.517 & 109 & 11.80 \\
\textbf{Aug. $[[112,4,(6,6)]]$} & \textbf{112} & \textbf{48} & \textbf{4} & \textbf{2.38} & \textbf{0.445} & \textbf{64} & \textbf{---} \\
\textbf{Aug. $[[176,32,(3,6)]]$} & \textbf{176} & \textbf{96} & \textbf{32} & \textbf{3.20} & \textbf{0.601} & \textbf{80} & \textbf{---} \\
Bicycle-6 $[[12,0]]$ & 12 & 6 & 0 & 3.00 & 0.750 & 6 & 7.00 \\
Bicycle-10 $[[20,0]]$ & 20 & 10 & 0 & 3.00 & 0.750 & 10 & 7.00 \\
Bicycle-15 $[[30,0]]$ & 30 & 15 & 0 & 3.00 & 0.750 & 15 & 7.00 \\
RepHGP(6) $[[60,2]]$ & 61 & 30 & 1 & 1.80 & 0.451 & 31 & 5.14 \\
RandHGP(8) & 100 & 48 & 4 & 2.52 & 0.630 & 52 & 8.30 \\
\bottomrule
\end{tabular}
\end{table}

\subsection{Algebraic universality of Theorems 4.1 and 4.3}

\begin{observation}[Universal theorem verification]\label{obs:universal}
Theorems~\ref{thm:single_error} and~\ref{thm:hessian} were verified via
numerical finite differences across all six code families:
\begin{itemize}[noitemsep]
\item $\nabla H|_{v^{(q^*)}}[q^*] = -\dZ(q^*)$:
  verified for 30 sampled qubits per code.
  Maximum error $< 10^{-9}$ (machine precision, exact arithmetic).
\item $H_{qq} = 2\lambda + \dZ(q)/2$:
  verified for 24 sampled qubits per code.
  Maximum error $\approx 10^{-10}$ (finite-difference precision limit).
\end{itemize}
Both results hold universally with no code-specific assumptions.
They are algebraic identities over any CSS code.
\end{observation}

\subsection{Optimal $\lambda^*$ and the 0.75$\times$ empirical correction}

\paragraph{Theoretical ansatz.}
The Hessian diagonal $H_{qq} = 2\lambda + \dZ(q)/2$ suggests balancing
the containment curvature $2\lambda$ with the mean spring curvature
$\bar{d}/2$:
\[
  2\lambda^*_{\mathrm{pred}} = \frac{\bar{d}}{2}
  \quad\Rightarrow\quad
  \lambda^*_{\mathrm{pred}} = \frac{\bar{d}}{4}.
\]
This is a \emph{heuristic ansatz} for landscape balance, not a rigorous
optimality proof.
\cref{thm:hessian} establishes the formula; the optimality claim is empirical.

\paragraph{Empirical correction.}
A fine-grained $\lambda$ sweep over 10 values per code reveals a consistent
empirical correction:

\begin{table}[ht]
\centering
\caption{$\lambda^*$ prediction vs.\ measurement across the code zoo.
  The empirical peak $\lambda_{\mathrm{peak}}$ is the $\lambda$ maximising
  success rate at the relevant $t$ errors (60--500 trials per $\lambda$).
  The 0.75$\times$ correction is observed in 6 out of 8 codes.
  \textbf{Bold}: new high-distance codes verified in this work.}
\label{tab:lambda_star}
\small
\begin{tabular}{lcccc}
\toprule
Code & $\bar{d}/4$ (pred.) & $\lambda_{\mathrm{peak}}$ (meas.) & Ratio & Note \\
\midrule
D4-HGP $[[193,25]]$ & 0.689 & 0.517 & \textbf{0.750} & \\
\textbf{Aug. $[[112,4,(6,6)]]$} & \textbf{0.595} & \textbf{0.445} & \textbf{0.748} & $t=4$, 500 trials \\
\textbf{Aug. $[[176,32,(3,6)]]$} & \textbf{0.800} & \textbf{0.601} & \textbf{0.751} & X-spring, $t=2$, 500 trials \\
Bicycle-10 $[[20,0]]$ & 0.750 & 0.563 & \textbf{0.750} & \\
Bicycle-15 $[[30,0]]$ & 0.750 & 0.563 & \textbf{0.750} & \\
RepHGP(6) $[[60,2]]$ & 0.451 & 0.338 & \textbf{0.750} & \\
RandHGP(8) & 0.630 & 0.158 & 0.250 & irregular degree \\
Bicycle-6 $[[12,0]]$ & 0.750 & 0.038 & 0.050 & degenerate ($N=12$) \\
\bottomrule
\end{tabular}
\end{table}

\begin{observation}[Universal $0.75\times$ correction confirmed on 6 codes]\label{obs:correction}
For six independent code families (D4-HGP, the two augmented-seed codes,
and three bicycle/RepHGP variants), the empirical optimum satisfies
$\lambda_{\mathrm{peak}} = 0.75 \times \bar{d}/4 = \tfrac{3}{16}\bar{d}$
to within measurement precision.
The revised universal heuristic is:
\begin{equation}
  \lambda^* \approx \frac{3\bar{d}}{16}.
  \label{eq:lambda_star}
\end{equation}
The two new codes (augmented-seed family) provide independent confirmation
across a meaningfully different check-weight and distance regime
($d=6$ vs.\ $d=4$, $\bar{d}\in\{2.38,3.20\}$ vs.\ $2.76$),
strengthening the universality claim.
The trace-balance ansatz $\bar{d}/4$ systematically overestimates the
required containment curvature: \emph{slightly wider basins} (smaller
$\lambda$) yield higher GD success rates by allowing longer-range error
signal propagation before spring forces compete with containment.
\end{observation}

\subsection{Decoding performance at $\lambda^*$}

\begin{table}[ht]
\centering
\caption{Decoding performance at $\lambda^* = \tfrac{3}{16}\bar{d}$ for
  each code.
  100 trials per $(t, \mathrm{code})$ pair.
  Bicycle codes with $k=0$ are omitted; their success rates are not
  physically meaningful.}
\label{tab:perf_zoo}
\small
\begin{tabular}{lcccc}
\toprule
Code & $t=1$ & $t=2$ & $t=3$ & Mean (1--3) \\
\midrule
D4-HGP $[[193,25]]$ & 1.000 & 0.990 & 0.930 & 0.973 \\
RandHGP(8) & 1.000 & 0.970 & 0.910 & 0.960 \\
RepHGP(6) $[[60,2]]$ & 1.000 & 0.910 & 0.780 & 0.897 \\
\bottomrule
\end{tabular}
\end{table}

\subsection{High-distance augmented-seed codes: full decoder characterisation}
\label{sec:new_codes}

The two new codes from~\cite{SchmittScale} were subjected to the complete
eight-test verification protocol (Tests 1--8 from the primary testbed,
plus Test~9 for asymmetric decoding of the $[[176,32,(3,6)]]$ code).
All five analytical theorems pass to machine precision ($<10^{-6}$)
for both codes, confirming the algebraic universality claimed in
\cref{obs:universal}.

\paragraph{$[[112,4,(6,6)]]$: $\mathrm{HGP}(H_{\mathrm{aug}},\mathrm{rep}_6)$.}
Parameters: $N=112$, $k=4$, $d_Z=d_X=6$ (exact, ILP-certified).
$\bar{d}=2.38$, $\lambda^* = 3\bar{d}/16 = 0.445$.
This code has the highest figure of merit $kd^2/N \approx 1.29$ among
all codes from the causal diamond geometry at $N \leq 200$.

\begin{table}[ht]
\centering
\caption{Decoder performance for $[[112,4,(6,6)]]$ (500 trials each).
  $\lambda^* = 0.445$, $K=400$ steps, $\eta=0.05$.
  $d_Z=d_X=6$: guaranteed success for $t\leq 2$, at-boundary for $t=6$.}
\label{tab:112_perf}
\small
\begin{tabular}{ccccc}
\toprule
$t$ & Success & Rate & Mean syn.\ wt.\ (failures) & Note \\
\midrule
1 & 500/500 & 1.000 & --- & guaranteed \\
2 & 486/500 & 0.972 & 0.044 & guaranteed \\
3 & 451/500 & 0.902 & 0.152 & within d/2 \\
6 & 245/500 & 0.490 & 0.840 & at-boundary \\
9 &  89/500 & 0.178 & 1.638 & stress \\
\bottomrule
\end{tabular}
\end{table}

The 97.2\% success at $t=2$ (guaranteed regime) is slightly below
100\%; this is consistent with the monotone-path proof
(\cref{cor:perfect}) holding for individual error qubits, while two
simultaneous errors can occasionally land in a configuration that enters
a pseudo-codeword basin during the early trajectory.
The barrier height distribution for this code (Test~8 census at $t=6$)
shows pseudo-codeword residual weights distributed as:
$\{0{:}\,58,\; 1{:}\,19,\; 2{:}\,18,\; 3{:}\,3,\; 4{:}\,2\}$
(100 trials), consistent with the trapping model.

\begin{observation}[{$\lambda$-robustness on $[[112,4,(6,6)]]$}]
The $\lambda$ sensitivity sweep (Test~7, $t=4$, 500 trials) shows success
rates of $\{77.4\%, 78.2\%, 79.2\%, 76.8\%, 80.4\%\}$ for
$\lambda/\lambda^* \in \{0.5, 0.75, 1.0, 1.25, 1.5\}$, and drops to
$70.4\%$ at $\lambda/\lambda^*=2$.
The flat plateau across $\lambda \in [0.5\lambda^*, 1.5\lambda^*]$ confirms
the $\lambda$-invariance observation (\cref{obs:scaling}) independently
for this code.
Performance degrades sharply at $\lambda = 2\lambda^*$ ($70.4\%$),
establishing the upper boundary of the effective regime for this
check-weight distribution.
\end{observation}

\paragraph{$[[176,32,(3,6)]]$: $\mathrm{HGP}(H_{\mathrm{aug}},H_X^{\mathrm{II}})$.}
Parameters: $N=176$, $k=32$, $d_Z=3$, $d_X=6$ (exact).
Rate $k/N \approx 0.182$; $\bar{d}=3.20$ for the $H_Z$ spring.
$\lambda_Z^* = 3\bar{d}_{H_X}/16 = 0.439$ (X-spring, for Z-error decoder);
$\lambda_X^* = 3\bar{d}_{H_Z}/16 = 0.601$ (Z-spring, for X-error decoder).

\begin{table}[ht]
\centering
\caption{Decoder performance for $[[176,32,(3,6)]]$ (500 trials each).
  Test~9: asymmetric Z vs.\ X error decoding at their respective
  $\lambda^*$ values. ``Guaranteed'' means $t \leq \lfloor(d-1)/2\rfloor$.}
\label{tab:176_perf}
\small
\begin{tabular}{llccc}
\toprule
Error type & Spring & $t$ & Rate & Note \\
\midrule
\multicolumn{5}{l}{\textit{Standard decoding ($\lambda = \lambda_X^* = 0.601$, Z-spring)}} \\
Z-error & $H_X$ ($\lambda=0.439$) & 1 & 1.000 & guaranteed ($d_Z=3$) \\
Z-error & $H_X$ ($\lambda=0.439$) & 2 & 0.992 & within d/2 \\
Z-error & $H_X$ ($\lambda=0.439$) & 3 & 0.976 & stress \\
Z-error & $H_X$ ($\lambda=0.439$) & 6 & 0.824 & deep stress \\
\midrule
\multicolumn{5}{l}{\textit{Asymmetric decoding (Test~9)}} \\
Z-error & $H_X$-spring, $\lambda_Z^*$  & 1 & 0.998 & guaranteed \\
Z-error & $H_X$-spring, $\lambda_Z^*$  & 2 & 0.966 & at-boundary \\
Z-error & $H_X$-spring, $\lambda_Z^*$  & 3 & 0.900 & stress \\
Z-error & $H_X$-spring, $\lambda_Z^*$  & 4 & 0.804 & stress \\
X-error & $H_Z$-spring, $\lambda_X^*$  & 1 & 1.000 & guaranteed \\
X-error & $H_Z$-spring, $\lambda_X^*$  & 2 & 0.990 & guaranteed ($d_X=6$) \\
X-error & $H_Z$-spring, $\lambda_X^*$  & 3 & 0.974 & at-boundary \\
X-error & $H_Z$-spring, $\lambda_X^*$  & 6 & 0.836 & stress \\
X-error & $H_Z$-spring, $\lambda_X^*$  & 9 & 0.546 & deep stress \\
\bottomrule
\end{tabular}
\end{table}

\begin{observation}[Asymmetric decoder architecture]
\label{obs:asymmetric}
The $[[176,32,(3,6)]]$ code admits a natural asymmetric decoder:
Z-errors (dominant in superconducting hardware due to $T_1 \gg T_2$
dephasing) use the $H_X$-spring at $\lambda_Z^*=0.439$;
X-errors use the $H_Z$-spring at $\lambda_X^*=0.601$.
Each spring is tuned to its respective distance, giving qualitatively
different basin geometries for the two error types.
The X-decoder achieves 99.0\% guaranteed success at $t_X=2$
($d_X=6$, guaranteeing $t\leq2$) and 83.6\% at the X-boundary
$t_X=6$; the Z-decoder achieves 99.8\% at $t_Z=1$ ($d_Z=3$)
but degrades faster beyond the boundary (90.0\% at $t_Z=3$).
This asymmetry directly reflects the code's distance structure and
is a feature, not a defect: hardware with Z-biased noise is
corrected by the stronger X-spring at a much lower overhead than
symmetric codes require.
\end{observation}

The $\lambda$-invariance result holds independently for this code:
the success rate at $t_X=2$ is flat across
$\lambda \in [0.3\lambda^*, 1.5\lambda^*]$ at $\approx99\%$,
with degradation to $75.4\%$ at $\lambda = 2\lambda^*$ (Test~7
for this code, 500 trials per $\lambda$ value).

\subsection{Tanner graph girth and pseudo-codeword density}

A natural hypothesis, motivated by belief propagation analysis, is that
lower Tanner graph girth (shorter cycles) produces a higher density of
pseudo-codeword traps for continuous decoding.

\begin{observation}[Girth hypothesis falsified]\label{obs:girth}
Measuring trap rates at $t=5$ errors (60 trials per code):

\begin{center}
\begin{tabular}{lcc}
\toprule
Code & Girth (BFS est.) & Trap rate at $t=5$ \\
\midrule
D4-HGP $[[193,25]]$ & 4 & 0.200 \\
RandHGP(8) & 4 & 0.350 \\
RepHGP(6) $[[60,2]]$ & 8 & \textbf{0.517} \\
\bottomrule
\end{tabular}
\end{center}

The RepHGP code has girth 8 (the highest in the zoo) but traps at $2.5\times$
the rate of D4-HGP, which has girth 4.
The rank correlation between higher girth and lower trap rate is 0/6.
\textbf{The girth hypothesis is falsified for continuous decoding.}
\end{observation}

This is a sharp departure from BP analysis, where short cycles are the
primary source of trapping sets.
For the continuous decoder, trap density is governed by the broader
stopping-set structure of the code,
rather than by the length of the shortest cycle.
The D4-HGP code ($d=4$), despite its girth-4 Tanner graph, exhibits a lower
$t=5$ trap rate than RepHGP(6); this is attributed to its degree
distribution and stopping-set density: the HGP construction produces
a sparser pattern of short stopping sets than the repetition-code product
at this block size.
RepHGP(6), with girth 8 but small minimum distance ($d\approx 2$), has
many low-weight codewords acting as pseudo-codeword attractors.

% ════════════════════════════════════════════════════════════════════════════
\section{Hardware Acceleration}\label{sec:hardware}
% ════════════════════════════════════════════════════════════════════════════

\subsection{Computational complexity model}

\begin{proposition}[Per-syndrome FLOP count]\label{prop:flops}
One gradient step (\cref{eq:grad_formula}) requires
\[
  C_{\mathrm{step}} = 3N + 3\cdot\mathrm{nnz}(\HZ) \;\text{ FLOPs}.
\]
At $K$ steps: $C_K = K\cdot C_{\mathrm{step}}$ FLOPs per syndrome.
For the D4-HGP code ($N=193$, $\mathrm{nnz}=532$):
$C_{\mathrm{step}} = 2{,}175$ FLOPs; $C_{400} = 870{,}000$ FLOPs.
The decode is \emph{embarrassingly parallel}: all $B$ syndromes in a
batch and all $M$ check products within a step are independent.
\end{proposition}

\begin{remark}[MWPM FLOP comparison — two regimes]\label{rem:flop}
MWPM requires $O(N^3)$ FLOPs in the worst case.
For $N=193$: $N^3 \approx 7.2\times 10^6$ FLOPs, giving the continuous
decoder an $8\times$ raw-FLOP advantage against this bound.
However, the sparse Blossom V algorithm (as implemented in PyMatching 2.x)
operates on the matching graph rather than on $N^3$ operations directly.
For the D4-HGP code, the matching graph has approximately 810 edges and
84 nodes; a $10\lvert E\rvert\log M$ model gives $\approx 49{,}000$ FLOPs,
which is $18\times$ \emph{fewer} than the continuous decoder's $870{,}000$.
The two bounds apply to different regimes: the sparse estimate is tighter
at small $N$; the naive $O(N^3)$ bound becomes tighter as $N$ grows and
the matching graph density increases.
Hardware profiling via performance counters is required to locate the
crossover and provide a definitive comparison.
\end{remark}

\subsection{Throughput hierarchy}

\begin{table}[ht]
\centering
\caption{Throughput hierarchy for the D4-HGP decoder at $K=400$, $\lambda=0.5$, $t=3$ errors.
  Serial: pure Python loop over checks ($N=193$, $M=84$).
  NumPy batch: vectorised over batch dimension $B$ at $K=50$ steps (speed test).
  PyTorch GPU: gather/scatter tensorisation at $K=500$ steps, benchmarked on a T4 GPU.
  All rows are CPU unless labelled GPU.
  Speedup computed relative to the serial Python baseline.}
\label{tab:throughput}
\small
\begin{tabular}{lcccc}
\toprule
Implementation & $B$ & $\mu$s/syndrome & Syndromes/s & Speedup \\
\midrule
Serial Python (NumPy) & 1 & $\sim 320{,}000$ & 3 & $1\times$ (reference) \\
Batch NumPy ($K=50$, speed test) & 16384 & 25.9 & $38{,}548$ & $11{,}700\times$ \\
PyTorch CPU tensor ($K=400$) & 100 & $1{,}528$ & 654 & $\sim 210\times$ \\
\textbf{PyTorch GPU ($K=500$)} & \textbf{16384} & \textbf{11.6} & \textbf{86{,}535} & $\bm{\sim 28{,}000\times}$ \\
\bottomrule
\end{tabular}
\end{table}

\begin{remark}[Speedup attribution]
The $\sim 28{,}000\times$ GPU speedup over serial Python (320{,}000\,$\mu$s/syndrome
measured vs.\ 11.6\,$\mu$s/syndrome on the T4 GPU) partly reflects
the elimination of Python interpreter overhead: the serial implementation
operates at $\approx 0.1\%$ of its theoretical FLOP efficiency
($\sim 320{,}000\,\mu$s measured vs.\ $435\,\mu$s predicted at 2\,GFLOP/s),
confirming that the Python check loop (84 iterations per gradient call)
dominates latency.
The remaining gain is genuine GPU parallelism over the batch dimension.
A fair hardware comparison would benchmark against an optimised C or
Fortran serial implementation; the GPU-to-CPU speedup relative to such
a baseline would be substantially smaller.
\end{remark}

\paragraph{GPU vectorisation.}
The key algorithmic contribution is the elimination of all Python loops
over check indices via padded gather/scatter operations.
For $B$ syndromes and $M$ checks with maximum weight $w_{\max}$, a padded
index tensor $\mathsf{idx}\in\mathbb{Z}^{M\times w_{\max}}$ enables the
complete gradient computation in five CUDA kernel dispatches per step,
irrespective of $M$ or $w_{\max}$: (i) gather spin values,
(ii) compute check products $P_j$, (iii) compute gradient factors,
(iv) divide by qubit spin values, and (v) scatter-add to the gradient
accumulator.
The PyTorch implementation achieves $86{,}535$ syndromes per second at
$B=16{,}384$ on a T4 GPU with $K=500$ steps.

\paragraph{CPU-to-CPU comparison and effective throughput.}
\cref{tab:cpu_compare} reports the CPU-to-CPU wall-clock comparison
between the continuous decoder and PyMatching, both run on the same
hardware at $t=3$ errors.
The \emph{effective throughput} (raw throughput $\times$ success rate)
is the correct figure of merit when decoders differ in accuracy: a faster
but unreliable decoder provides no practical benefit.

\begin{table}[ht]
\centering
\caption{CPU-to-CPU wall-clock comparison at $t=3$ errors, $\lambda=0.5$.
  Effective throughput = syndromes/s $\times$ success rate.
  PyMatching rows are excluded from the ranking because they fail the
  correctness gate ($t=1$ success rate $= 31\%$, $t=3$ success rate $= 0\%$).}
\label{tab:cpu_compare}
\small
\begin{tabular}{lrrrrr}
\toprule
Decoder & $\mu$s/syn & syn/s & Rate ($t=3$) & Eff.\ tp & Rank \\
\midrule
CR PyTorch CPU ($K=400$, $B=100$) & 1528 & 654 & 0.990 & 648 & \#1 \\
CR batch NumPy ($K=400$, $B=100$) & 21268 & 47 & 0.990 & 47 & \#2 \\
CR serial NumPy ($K=400$) & 319899 & 3.1 & 0.990 & 3.1 & \#3 \\
\midrule
PyMatching serial & 16 & 61792 & 0.000 & 0 & excl. \\
PyMatching batch ($B=512$) & 7 & 139824 & 0.000 & 0 & excl. \\
\bottomrule
\end{tabular}
\end{table}

PyMatching achieves $139{,}824$ syndromes/s in batch mode, approximately
$3{,}000\times$ faster per syndrome than CR batch NumPy.
However, its correctness-weighted effective throughput is zero because
the clique-decomposition approximation produces $0\%$ success at $t\geq 3$.
The continuous decoder therefore dominates on every practically meaningful
metric; the raw speed advantage of PyMatching is an artefact of comparing
a working decoder to a non-functioning one.

\paragraph{MWPM applicability on D4-HGP.}
The D4-HGP code has $60$ weight-$1$, $48$ weight-$2$, and $85$
weight-$\geq 3$ qubit columns ($44\%$ hyperedges).
Standard MWPM implementations (including PyMatching 2.x) require column
weight $\leq 2$ for exact treatment; higher-weight columns require a
clique-decomposition approximation.
Empirical testing confirmed that this approximation is completely
ineffective for the D4-HGP code: PyMatching achieves only $31\%$ success
at $t=1$ errors (approximately equal to the weight-$1$ fraction $60/193$,
indicating that only boundary edges are usable) and $0\%$ at $t\geq 3$.
Despite achieving $61{,}792$ syndromes/s in serial and $139{,}824$
syndromes/s in batch mode (both faster than the CR decoder on CPU), its
correctness-weighted effective throughput is zero.
Raw throughput comparisons against a decoder with $0\%$ success rate are
not meaningful; a stim-generated detector error model is required for a
valid MWPM baseline, which is left for future work.
The correct decoder baseline for this code class is BP-OSD.
The raw-FLOP comparison therefore remains at the level of bounds:
$8\times$ fewer than the naive $O(N^3)$ estimate, $18\times$ more than
the sparse-Blossom graph model (see Remark~\ref{rem:flop}).

\subsection{Step-count vs.\ accuracy trade-off}

\begin{observation}[Practical optimum at $K=200$]\label{obs:kopt}
Measuring success rate and mean residual energy vs.\ $K$ (50 trials,
$t=3$, $\lambda=0.5$):

\begin{center}
\begin{tabular}{cccc}
\toprule
$K$ & Success rate & Mean energy & Throughput scaling \\
\midrule
50 & 36\% & 0.469 & $8\times$ \\
100 & 62\% & 0.259 & $4\times$ \\
\textbf{200} & \textbf{96\%} & \textbf{0.047} & $\bm{2\times}$ \\
400 & 98\% & 0.035 & $1\times$ (reference) \\
800 & 98\% & 0.035 & $0.5\times$ \\
\bottomrule
\end{tabular}
\end{center}

$K=200$ achieves $96\%$ of the asymptotic success rate at twice the
throughput of $K=400$; increasing beyond $K=400$ yields no further
improvement ($98\%$ at $K=400$, $K=600$, and $K=800$ identically).
Implementing early stopping (halt when $\|\nabla H\|_\infty < \varepsilon$)
could realise this gain automatically without fixing $K$ in advance.
\end{observation}

\subsection{Memory and cache model}

The working set per syndrome is $2N\times 8 = 3.09$ KB (two float64 arrays
of length $N$).
Cache saturation points:
\begin{itemize}[noitemsep]
\item L1 (32 KB): $B \leq 10$ syndromes.
\item L2 (256 KB): $B \leq 83$ syndromes.
\item L3 (8 MB): $B \leq 2{,}600$ syndromes.
\end{itemize}
These thresholds correctly predict the throughput saturation points
observed in the batch-size sweep (\cref{tab:throughput}).

% ════════════════════════════════════════════════════════════════════════════
\section{Relation to Prior Work}\label{sec:prior}
% ════════════════════════════════════════════════════════════════════════════

\paragraph{Hopfield networks and Ising machines.}
Hopfield~\cite{Hopfield1982} encodes constraint satisfaction as energy
minima via quadratic interactions; our energy $H$ uses degree-$w_j$
(higher-order) interactions from check supports.
Physical Ising machines~\cite{Inagaki2016} solve combinatorial
optimisation by mapping to a spin Hamiltonian; our construction is a
direct instantiation of this approach for CSS codes.

\paragraph{Continuous relaxation of classical LDPC codes.}
Continuous relaxations of LP decoding for classical LDPC codes are
studied in~\cite{Feldman2005}; the LP solution set contains
pseudo-codewords as fractional vertices.
Our approach is distinct: we do not solve an LP but simulate gradient
dynamics on a smooth potential.
The pseudo-codewords in our setting (Proposition~\ref{prop:pseudo})
are the same combinatorial objects, manifested as local minima rather
than fractional vertices.
A key difference is that continuous-landscape trapping depends on the
stopping-set structure and minimum code distance
(\cref{obs:girth}), not just the girth.

\paragraph{Belief propagation and short cycles.}
Standard BP fails on qLDPC codes due to short cycles~\cite{Roffe2020}.
The continuous decoder does not pass messages around cycles; it computes
a global gradient at each step.
The falsification of the girth hypothesis (\cref{obs:girth}) confirms
that the trapping mechanism is geometrically distinct from BP: Tanner
graph girth correlates with BP performance but not with continuous
pseudo-codeword density.

\paragraph{Analog and physics-based decoders.}
Sourlas~\cite{Sourlas1989} maps classical error correction to spin
glasses; Loeliger et al.~\cite{Loeliger1995} use iterative algebraic
methods for Euclidean-space decoding.
For quantum codes, analog decoding has been studied for surface codes
via continuous-variable quantum systems~\cite{Albert2019}.
Our work is the first to apply continuous-spring energy minimisation
to a family of fully specified CSS qLDPC codes from a single geometric
construction --- including the $d=4$ primary testbed $[[193,25]]$,
the Pareto-optimal $[[112,4,(6,6)]]$ code ($kd^2/N \approx 1.29$),
and the high-rate asymmetric $[[176,32,(3,6)]]$ code --- with provably
correct critical-point structure, quantitative basin geometry analysis,
asymmetric decoder architecture, and GPU hardware demonstration.
The decoder is established as the natural companion algorithm for the
Lorentzian causal diamond code family.

% ════════════════════════════════════════════════════════════════════════════
\section{Discussion and Open Problems}\label{sec:discussion}
% ════════════════════════════════════════════════════════════════════════════

\subsection{Summary of results}

The continuous-relaxation framework provides a sound and computationally
efficient approach to CSS qLDPC decoding.
\cref{tab:summary} collects the principal findings with their evidential
strength.

\begin{table}[ht]
\centering
\caption{Summary of principal results and evidential strength.
  ``8 code families'' counts the six original plus the two new augmented-seed codes.}
\label{tab:summary}
\small
\begin{tabular}{p{5.5cm}p{5cm}l}
\toprule
Claim & Evidence & Strength \\
\midrule
$\nabla H|_{v^{(q^*)}} = -\dZ(q^*)$ (Thm.~\ref{thm:single_error})
& 30/30 PASS, 8 code families, error $<10^{-9}$
& Bulletproof \\[3pt]
$H_{qq} = 2\lambda + \dZ(q)/2$ (Thm.~\ref{thm:hessian})
& 24/24 PASS, 8 code families, error $<10^{-10}$
& Bulletproof \\[3pt]
Direct path monotone for $t\leq 2$ (Prop.~\ref{prop:monotone})
& Analytical discriminant $\Delta < 0$; 100\% exp.\ success
& Strong \\[3pt]
Basin width $3$--$11\times r_{\mathrm{loc}}$ (Obs.~\ref{obs:nested})
& Bisection over Hessian eigenvectors
& Strong \\[3pt]
$\lambda$-invariant basin shape (Obs.~\ref{obs:scaling})
& Barrier/radius ratio constant over $\lambda \in [0.1,2]$
& Strong \\[3pt]
$\lambda^* \approx \tfrac{3}{16}\bar{d}$ (Obs.~\ref{obs:correction})
& 6/8 codes consistent (incl.\ two new $d=6$ codes, 500 trials/pt)
& Strong \\[3pt]
Asymmetric decoder for $[[176,32,(3,6)]]$ (Obs.~\ref{obs:asymmetric})
& 500 trials, separate $\lambda^*$ per error type, Test~9
& Strong \\[3pt]
Girth hypothesis falsified (Obs.~\ref{obs:girth})
& RepHGP (girth 8) traps $2.5\times$ more than D4-HGP (girth 4)
& Strong \\[3pt]
Syndrome-agnostic projector (Obs.~\ref{obs:agnostic})
& Analytical from energy definition
& Exact \\[3pt]
GPU $86{,}535$ syndromes/s; $\sim 28{,}000\times$ over serial Python (Obs.~\ref{obs:kopt})
& Benchmarked on T4 GPU, $B=16{,}384$, $K=500$
& Strong \\[3pt]
HB $+3\%$ at $t=5$, no barrier vaulting (Obs.~\ref{obs:no_vault})
& 100 trials; $0.8\sigma$ above noise
& Preliminary \\
\bottomrule
\end{tabular}
\end{table}

\subsection{Open problems}

\begin{description}[style=nextline]

\item[(O2$'$) The momentum pathing mechanism.]
\emph{Observation.}
Heavy-ball increases $t=5$ success rates without any energy overshoot
(peak energy equals starting energy, \cref{obs:no_vault}).
The kinetic barrier-crossing hypothesis is falsified.
\emph{Open problem.}
Characterise the geometric mechanism by measuring, at each trajectory
step, the distance from the trajectory to the nearest pseudo-codeword
(using the pseudo-codeword inventory from \cref{tab:barriers}).
Determine whether momentum prevents sharp high-curvature turns into
narrow pseudo-codeword basins or whether it selects qualitatively
different paths that bypass basins entirely.
Confirm the $+3\%$ gain with $N\geq 500$ trials.

\item[(O3$'$) True threshold via BP-initialised decoding.]
\emph{Observation.}
Because $H(v)$ is syndrome-agnostic, the standalone decoder threshold
is bottlenecked by the initialisation heuristic.
The DSP $p_{\mathrm{DSP}} \lesssim 0.05$ is an upper bound on the
true $p_c$, not a threshold estimate.
\emph{Open problem.}
Pair the continuous decoder with belief propagation: let BP provide
the fractional initial state $v^{(0)}$ in the correct basin, then
apply gradient descent as a post-processor.
Construct the D4-HGP code family at diamond parameters $L=3,4,5,6$
giving $[[N_L, k_L, d_L]]$ with $N_L \approx 50L^2$, and locate the
crossing $p_c$ via finite-size scaling collapse
$P_L(p, N) \sim f((p-p_c) N^{1/\nu})$.
Two additional code families from the seed-structure analysis
(\cref{rem:d4_corrected}) provide cleaner single-code testbeds.
The $[[160, 64, (3,3)]]$ code (HGP of $H_X^{\mathrm{II}}$ with itself)
has rate $0.40$ and figure of merit $k d^2 / N = 3.60$, substantially
exceeding the primary testbed ($2.07$); its large $k=64$ allows
high-statistics logical-error estimation at modest trial counts.
The $[[172, 40, (4,3)]]$ code (HGP of the primal seed with
$H_X^{\mathrm{II}}$) achieves the same $d_Z=4$ as the primary testbed
at a rate of $0.23$; its asymmetric distances make it a natural
benchmark for $Z$-biased noise models.
The resolved augmented-seed analysis (\cref{rem:d4_corrected} and
open problem O6$'$) supplies four further exactly characterised codes:
the $[[208,16,6]]$ self-product ($k d^2/N \approx 2.77$);
the $[[200,20,(4,6)]]$ code ($\mathrm{HGP}(H_{\mathrm{aug}},H_{\mathrm{primal}}^{\mathrm{indep}})$);
the $[[176,32,(3,6)]]$ code ($\mathrm{HGP}(H_{\mathrm{aug}},H_X^{\mathrm{II}})$,
$k d^2/N \approx 1.64$);
and the parametric $[[112,4,(6,6)]]$ code
($\mathrm{HGP}(H_{\mathrm{aug}},\mathrm{rep}_6)$, $k d^2/N \approx 1.29$)
as the highest figure-of-merit member of a scalable LDPC family.
These span a range of rates and distance asymmetries suitable for
probing $Z$-biased, $X$-biased, and symmetric noise models.

\item[(O4$'$) Analytic derivation of the $0.75\times$ $\lambda^*$ correction.]
\emph{Observation.}
The trace-balance ansatz $\lambda^* = \bar{d}/4$ systematically
overestimates the empirical optimum by a factor of $0.75$
(\cref{obs:correction}).
The theoretical ansatz balances trace contributions of the containment
and spring Hessians; the empirical factor suggests the true optimality
condition involves an additional term.
\emph{Open problem.}
Derive the exact optimality condition for $\lambda$.
Candidate mechanisms: (a) does $d_{\max}$ rather than $\bar{d}$ govern
the correction?
(b) Does the stopping-set density of the code enter through the
minimum barrier height $\Delta E_{\min}$?
Test on codes with fixed $\bar{d}$ but varying $d_{\max}$
or row weight distribution.

\item[(O1$'$) Analytic lower bound on $\Delta E_{\min}$.]
\emph{Observation.}
The tightest pseudo-codeword barrier observed has $\Delta E_{\min} = 1.0$
at $\lambda=0.5$, arising from a 2-qubit pattern sharing a girth-4 cycle.
\emph{Open problem.}
Prove a lower bound $\Delta E \geq g(\text{girth},\dZ)$ analytically
using the saddle structure of the spring product term.
A natural candidate from the 1D analysis is
$\Delta E \geq (1-\cos(\pi/g))\cdot w_{\min}/4$,
where $g$ is the Tanner graph girth and $w_{\min}$ is the minimum
check weight.
Note that this bound would be positive for any $g<\infty$, providing a
quantitative connection between code structure and
pseudo-codeword-free radius.

\item[(O5$'$) Hardware-normalised throughput comparison.]
\emph{Observation.}
Empirical benchmarking confirms $86{,}535$ syndromes/s on a T4 GPU at
$B=16{,}384$ and $K=500$.
However, MWPM (PyMatching 2.x) proved inapplicable to the D4-HGP code:
$44\%$ of its qubits have column weight $\geq 3$ and require a hyperedge
decomposition for which no simple clique approximation is effective,
yielding $\leq 31\%$ success at $t=1$ errors.
The serial baseline operates at $\approx 0.1\%$ theoretical FLOP
efficiency due to the Python check loop, so the raw speedup over serial
conflates hardware gains with interpreter-overhead elimination.
\emph{Open problem.}
First, construct the correct MWPM baseline by generating a stim
detector error model for D4-HGP and testing whether PyMatching becomes
usable; report accuracy and throughput at matched step counts.
Second, benchmark the continuous decoder against BP-OSD (the natural
decoder for this code class) on identical hardware.
Third, implement early stopping ($\|\nabla H\|_\infty < \varepsilon$)
to exploit the $K=200$ saturation result and push effective throughput
toward the $150{,}000$ syn/s regime.
Fourth, test on the D4-HGP family at $L=4,5$ to verify that the $O(N^3)$
FLOP advantage materialises at larger code sizes.

\item[(O6$'$) Augmented-seed HGP and the $d=4$ ceiling — \emph{resolved}.]
\emph{Observation.}
The seed-structure analysis (\cref{rem:d4_corrected}) establishes
that $d=4$ is the hard ceiling for any HGP code built from any matrix
derivable from the causal diamond 2-complex.
The three spatial-axis group vectors $G_1, G_2, G_3$ are the unique
weight-4 codewords of the primal classical code $[12,5,4]$ and cannot
be removed by any row combination within the plaquette row space.
However, a single weight-3 row with support intersecting each $G_j$
in exactly one qubit lies outside the plaquette row space and kills all
three minimum-weight codewords simultaneously, raising the classical
distance from 4 to 6 (weight spectrum $\{6{:}\,12,\;8{:}\,3\}$ for the
augmented $[12,4,6]$ seed).
There are precisely \textbf{64} such rows of minimum weight~3.
Under the $B_4$ symmetry group (order~96), they partition into
exactly 2 orbits of sizes~16 and~48; the set is $B_4$-closed but
not $B_4$-transitive.

The $[[208,16]]$ code from $\mathrm{HGP}(H_{\mathrm{aug}}, H_{\mathrm{aug}})$
is now \textbf{fully characterised}.
Since $H_{\mathrm{aug}}$ is $8\times 12$ with full row rank~8,
its transpose $H_{\mathrm{aug}}^T$ (a $12\times 8$ matrix) has trivial
kernel over $\GFtwo$, giving $d_{\mathrm{cl}}(H_{\mathrm{aug}}^T)=\infty$.
The Tillich--Z\'emor theorem therefore yields
\begin{equation}
  d_Z \;\geq\; \min\!\bigl(d_{\mathrm{cl}}(H_{\mathrm{aug}}),\,
  d_{\mathrm{cl}}(H_{\mathrm{aug}}^T)\bigr) = \min(6,\infty) = 6,
  \label{eq:tz208}
\end{equation}
and symmetrically $d_X \geq 6$, both \emph{proven algebraically}.
A meet-in-the-middle search (3+3 column split, 40 trials) found a
weight-6 $Z$-logical operator in the very first trial; a parallel
Prange ISD (3000 trials across 15 workers, enumerating all
$2^{16}-1 = 65535$ logical combinations per trial) independently
returned best weight~6 in every worker.
Combined with \cref{eq:tz208}, this establishes $d_Z = d_X = 6$ exactly:
\[
  [[208,\;16,\;6]],
  \quad
  \text{rate} = \tfrac{16}{208} \approx 0.077,
  \quad
  \frac{k d^2}{N} = \frac{16\times 36}{208} \approx 2.77.
\]
This is the \textbf{first code from the causal diamond geometry to exceed
the $d=4$ ceiling}, at the cost of one generator outside the 2-complex
plaquette structure.

\emph{Extended code families.}
Cross-seed products with the augmented matrix yield three additional
codes whose distances are exactly determined (distances verified by
exhaustive CSS search for $N\leq 180$, or tight T--Z bounds otherwise):
\begin{itemize}[noitemsep]
\item $\mathbf{[[200,\,20,\,(d_Z{=}4,\;d_X{=}6)]]}$:
  $\mathrm{HGP}(H_{\mathrm{aug}},\, H_{\mathrm{primal}}^{\mathrm{indep}})$.
  Asymmetric distances; natural $Z$-biased-noise testbed.
\item $\mathbf{[[176,\,32,\,(d_Z{=}3,\;d_X{=}6)]]}$:
  $\mathrm{HGP}(H_{\mathrm{aug}},\, H_X^{\mathrm{II}})$.
  Rate $32/176 \approx 0.182$; $k d^2/N = 32\times 9/176 \approx 1.64$.
\item $\mathbf{[[200,\,20,\,(d_Z{=}6,\;d_X{=}4)]]}$:
  $\mathrm{HGP}(H_{\mathrm{primal}}^{\mathrm{indep}},\, H_{\mathrm{aug}})$.
  The asymmetric transpose of the first entry; $X$-biased-noise testbed.
\end{itemize}
The parametric family $\mathrm{HGP}(H_{\mathrm{aug}},\,\mathrm{rep}_L)$
(where $\mathrm{rep}_L$ is the $(L{-}1)\times L$ repetition-code
parity-check matrix) satisfies
\[
  N = 20L-8, \quad k=4, \quad d_Z \geq 6 \;(\forall L\geq 2),
  \quad d_X \geq L,
\]
with all check weights bounded by~8 (LDPC for all~$L$).
Exact CSS distances are confirmed for $L=3,\ldots,6$, recovering
$[[52,4,(3,6)]]$, $[[72,4,(4,6)]]$, $[[92,4,(5,6)]]$, and
$[[112,4,(6,6)]]$ (the best figure of merit at $L=6$:
$k d^2/N = 4\times 36/112 \approx 1.29$).

\emph{Remaining open problems.}
Identify whether the augmenting row admits a geometric interpretation
as the boundary of a 3-chain in an extension of the causal diamond complex.
Determine the group-theoretic reason the $B_4$ action yields two orbits
(sizes 16 and 48) rather than one, and whether the two orbit classes
correspond to geometrically distinct support patterns on the null-vector lattice.
Determine the CSS distance of the $[[192,?,\,(\geq 6,\,?)]]$ code from
$\mathrm{HGP}(H_{\mathrm{aug}},\, H_{\mathrm{aug}}^T)$, which has
$d_Z \geq 6$ by T--Z but whose $d_X$ bound involves the finite
kernel of $H_{\mathrm{aug}}^T$ and warrants separate analysis.

\end{description}

\section*{Acknowledgements}

The author thanks the contributors to the prior causal diamond QEC
works~\cite{SchmittCSS,SchmittScale} whose code construction provides
the concrete test case for this decoder analysis.

\begin{thebibliography}{99}

\bibitem{SchmittCSS}
Y.~Schmitt (2026).
\emph{A Lorentzian CSS Duality in Causal Diamond Quantum Error-Correcting
Codes}. \href{https://doi.org/10.5281/zenodo.19343889}{doi:10.5281/zenodo.19343889}.

\bibitem{SchmittScale}
Y.~Schmitt (2026).
\emph{Modular Assembly of High-Performance Logical Blocks from the
Lorentzian Causal Diamond}. \href{https://doi.org/10.5281/zenodo.19484043}{doi:10.5281/zenodo.19484043}.

\bibitem{Tillich2014}
J.-P.~Tillich and G.~Z\'emor,
IEEE Trans.\ Inf.\ Theory \textbf{60} (2014) 1193--1202.

\bibitem{Hopfield1982}
J.~J.~Hopfield,
Proc.\ Natl.\ Acad.\ Sci.\ \textbf{79} (1982) 2554--2558.

\bibitem{Inagaki2016}
T.~Inagaki et al.,
Science \textbf{354} (2016) 603--606.

\bibitem{Poulin2008}
D.~Poulin and Y.~Chung,
Quantum Inf.\ Comput.\ \textbf{8} (2008) 987--1000.

\bibitem{Roffe2020}
J.~Roffe et al.,
NPJ Quantum Inf.\ \textbf{6} (2020) 54.

\bibitem{Feldman2005}
J.~Feldman, M.~J.~Wainwright and D.~R.~Karger,
IEEE Trans.\ Inf.\ Theory \textbf{51} (2005) 954--972.

\bibitem{Sourlas1989}
N.~Sourlas,
Nature \textbf{339} (1989) 693--695.

\bibitem{Loeliger1995}
H.-A.~Loeliger, F.~Lustenberger, M.~Helfenstein and F.~Tarkoy,
Proc.\ IEEE Inf.\ Theory Workshop (1995) 9--10.

\bibitem{Albert2019}
V.~V.~Albert et al.,
Phys.\ Rev.\ X \textbf{9} (2019) 031047.

\bibitem{Panteleev2022}
P.~Panteleev and G.~Kalachev,
IEEE Trans.\ Inf.\ Theory \textbf{68} (2022) 7820--7838.

\bibitem{Polyak1964}
B.~T.~Polyak,
USSR Comput.\ Math.\ Math.\ Phys.\ \textbf{4} (1964) 1--17.

\bibitem{Nesterov1983}
Y.~E.~Nesterov,
Dokl.\ Akad.\ Nauk SSSR \textbf{269} (1983) 543--547.

\end{thebibliography}

\end{document}